% !TEX root = 抽象代数前两章全部习题详解.tex
\clearpage\section{第2章：群}

\subsection{习题1：\texorpdfstring{阶互素时落入正规子群}{阶互素时落入正规子群}}
\label{ex:2:1}

\paragraph{题目}
有限群$G$中，$N\trianglelefteq G$且$\gcd(|N|,|G/N|)=1$。若$a$的阶整除$|N|$，则$a\in N$。

\paragraph{考点与思路}
同时利用元素在原群和商群中的阶，以及两个阶数的互素性。

\paragraph{所用结论与原文依据}
\noteref{order}；\noteref{lagrange}；\noteref{quotient}。

% 原解来源：chapters/Part 2/chapter2-exercises.tex:7；第2章原\exref{2}{1}。

\begin{proof}
设$m=\ord(a)$、$d=\ord(aN)$。由$m\mid|N|$，$m$有限；自然同态保持幂，故
\[
(aN)^m=a^mN=N.
\]
所以$d\mid m$，从而$d\mid|N|$。另一方面，$aN$是有限群$G/N$的元素，由\proofref{lagrange}{拉格朗日定理}，$d\mid|G/N|$。

因此$d$是$|N|$和$|G/N|$的公因子。题设两数互素，故$d=1$，即$aN=N$。陪集相等的判据于是给出$a\in N$。
\end{proof}

\paragraph{复习提示}
商群的单位元是陪集$N$，不是原群元素$e$。


\subsection{习题2：\texorpdfstring{互素阶的分解}{互素阶的分解}}
\label{ex:2:2}

\paragraph{题目}
若$\operatorname{ord}(c)=rs$、$\gcd(r,s)=1$，则$c=ab$，其中$a,b\in\langle c\rangle$的阶分别为$r,s$。

\paragraph{考点与思路}
先用贝祖等式分解指数1，再使用幂的阶公式。

\paragraph{所用结论与原文依据}
\noteref{powerorder}；\noteref{powers}。补充前置工具：整数贝祖等式。

% 原解来源：chapters/Part 2/chapter2-exercises.tex:19；第2章原\exref{2}{2}。

\begin{proof}
由$\gcd(r,s)=1$，存在整数$u,v$使$us+vr=1$。定义$a=c^{us}$、$b=c^{vr}$。显然$a,b\in\langle c\rangle$，并且
\[
ab=c^{us+vr}=c.
\]
由$us\equiv1\pmod r$知$\gcd(u,r)=1$；由$vr\equiv1\pmod s$知$\gcd(v,s)=1$。因而
\begin{align*}
\gcd(rs,us)&=s\gcd(r,u)=s,\\
\gcd(rs,vr)&=r\gcd(s,v)=r.
\end{align*}
按原笔记的\proofref{powerorder}{幂的阶公式}，
\[
\ord(a)=\frac{rs}{\gcd(rs,us)}=r,\qquad
\ord(b)=\frac{rs}{\gcd(rs,vr)}=s.
\]
这就得到题设要求的分解；$r=1$或$s=1$时同一计算仍适用。
\end{proof}

\paragraph{复习提示}
指数$u,v$可以为负数；群中的整数幂已经有定义。


\subsection{习题3：\texorpdfstring{互素指数的开方}{互素指数的开方}}
\label{ex:2:3}

\paragraph{题目}
若$\operatorname{ord}(a)=m$且$\gcd(k,m)=1$，则$x^k=a$在$\langle a\rangle$内有唯一解。

\paragraph{考点与思路}
把循环群中的幂方程化成指数的同余方程。

\paragraph{所用结论与原文依据}
\noteref{order}；\noteref{powers}。补充前置工具：整数贝祖等式与互素整除。

% 原解来源：chapters/Part 2/chapter2-exercises.tex:39；第2章原\exref{2}{3}。

\begin{proof}
由贝祖等式，存在整数$u,v$使$ku+mv=1$。取$x=a^u$，则$x\in\langle a\rangle$，且
\[
x^k=a^{ku}=a^{1-mv}=a(a^m)^{-v}=a.
\]
所以解存在。

若$y\in\langle a\rangle$也满足$y^k=a$，可写$y=a^t$。则$a^{kt-1}=e$，由$\ord(a)=m$得$m\mid kt-1$。另一方面$m\mid ku-1$，相减得$m\mid k(t-u)$。因为$\gcd(k,m)=1$，故$m\mid t-u$，于是$y=a^t=a^u=x$。因此解唯一。
\end{proof}

\paragraph{复习提示}
唯一性仅在题设指定的循环子群内讨论，不可直接扩大到整个群。


\subsection{习题4：\texorpdfstring{$ab$与$ba$同阶}{ab与ba同阶}}
\label{ex:2:4}

\paragraph{题目}
设$a,b$为群$G$的任意两个元素，证明$ab$与$ba$具有相同的阶，包括无限阶的情形。

\paragraph{考点与思路}
沿原笔记用共轭联系两个乘积，比较所有使幂等于单位元的正指数。

\paragraph{所用结论与原文依据}
\noteref{order}；\noteref{glaws}。

% 原解来源：chapters/Part 2/chapter2-exercises.tex:57；第2章原\exref{2}{4}。

\begin{proof}
由结合律，$a^{-1}(ab)a=ba$。对每个正整数$m$，相邻的$aa^{-1}$消去，得到
\[
(ba)^m=(a^{-1}(ab)a)^m=a^{-1}(ab)^ma.
\]
若$(ab)^m=e$，右边等于$e$，所以$(ba)^m=e$；反之，若$(ba)^m=e$，在等式两边左乘$a$、右乘$a^{-1}$，得到$(ab)^m=e$。

因此两者满足正幂等于单位元的正整数集合相同。该集合非空时，最小正整数相同，即有限阶相同；该集合为空时，两者都为无限阶。
\end{proof}

\paragraph{复习提示}
只比较有限阶还不够，要交代一个乘积无限阶时另一个也无限阶。


\subsection{习题5：\texorpdfstring{给定阶元素的奇偶性}{给定阶元素的奇偶性}}
\label{ex:2:5}

\paragraph{题目}
有限群中，阶恰为$n>2$的元素个数为偶数。

\paragraph{考点与思路}
在所有$n$阶元素组成的集合中按取逆配对，利用$n>2$排除不动点。

\paragraph{所用结论与原文依据}
\noteref{order}；\noteref{glaws}。

% 原解来源：chapters/Part 2/chapter2-exercises.tex:73；第2章原\exref{2}{5}。

\begin{proof}
令$X=\{a\in G:\ord(a)=n\}$。若$\ord(a)=n$，则$(a^{-1})^m=e$当且仅当$a^m=e$，故$\ord(a^{-1})=n$。所以取逆把$X$映入自身，且连续取两次逆得到原元素。

若$a=a^{-1}$，则$a^2=e$，所以$\ord(a)\mid2$，这与$n>2$矛盾。因此$X$中的元素按$\{a,a^{-1}\}$分成不相交的二元集合。$G$有限，所以$X$有限，其元素个数为偶数；$X$为空时其元素个数为0，也符合结论。
\end{proof}

\paragraph{复习提示}
需要配对的集合是阶恰为$n$的元素，不能把阶整除$n$混在一起。


\subsection{习题6：\texorpdfstring{对称群的中心}{对称群的中心}}
\label{ex:2:6}

\paragraph{题目}
$Z(S_n)=\{e\}$，$n>2$。

\paragraph{考点与思路}
中心元素与每个对换交换，因而保持每个二元子集。

\paragraph{所用结论与原文依据}
\noteref{center}；\noteref{conjugation}。

% 原解来源：chapters/Part 2/chapter2-exercises.tex:88；第2章原\exref{2}{6}。

\begin{proof}
设$\sigma\in Z(S_n)$。对任意不同的$i,j$，由\proofref{conjugation}{共轭公式}及中心的定义，
\[
(\sigma(i)\ \sigma(j))=\sigma(ij)\sigma^{-1}=(ij).
\]
两个对换相同当且仅当它们交换的是同一个二元集合，因此
\[
\{\sigma(i),\sigma(j)\}=\{i,j\}.
\]
固定任意$i$。由于$n>2$，可选择两个与$i$不同且彼此不同的$j,k$。上式分别用于$\{i,j\}$及$\{i,k\}$，得到
\[
\sigma(i)\in\{i,j\}\cap\{i,k\}=\{i\}.
\]
所以$\sigma(i)=i$。$i$任意，故$\sigma=e$。反之恒等置换属于中心，因而$Z(S_n)=\{e\}$。
\end{proof}

\paragraph{复习提示}
条件$n>2$保证可以选取第三个点；$S_2$本身交换，中心并不平凡。


\subsection{习题7：\texorpdfstring{有理数加法群的有限生成子群}{有理数加法群的有限生成子群}}
\label{ex:2:7}

\paragraph{题目}
证明$(\Q,+)$的任一有限生成子群都是循环群。

\paragraph{考点与思路}
把有限多个有理数通分，将生成子群转化为整数的线性组合，再用最大公因数和贝祖等式确定生成元。

\paragraph{所用结论与原文依据}
\noteref{generated}；\noteref{cyclic}。补充前置工具：整数贝祖等式。

% 原解来源：chapters/Part 2/chapter2-exercises.tex:105；第2章原\exref{2}{7}。

\begin{proof}
设$H=\langle q_1,\ldots,q_t\rangle$，这里加法群中的生成意味着
\[
H=\left\{\sum_{i=1}^t k_iq_i:k_i\in\Z\right\}.
\]
若生成集为空或全部$q_i=0$，则$H=\{0\}=\langle0\rangle$，已经循环。以下设至少一个$q_i\ne0$。

取所有$q_i$的正公分母$D$，写$q_i=m_i/D$，其中$m_i\in\Z$。令$d=\gcd(m_1,\ldots,m_t)>0$。因$d$整除每个$m_i$，对任意$k_i\in\Z$都有
\[
\sum_i k_iq_i=\frac{\sum_i k_im_i}{D}\in\frac dD\Z.
\]
因此$H\subseteq(d/D)\Z$。

反过来，由贝祖等式，存在整数$u_1,\ldots,u_t$使$\sum_i u_im_i=d$。除以$D$得
\[
\frac dD=\sum_i u_iq_i\in H.
\]
由于$H$是加法子群，它包含$d/D$的全部整数倍，故$(d/D)\Z\subseteq H$。综上，
\[
H=\frac dD\Z=\left\langle\frac dD\right\rangle.
\]
\end{proof}

\paragraph{复习提示}
“每个有限生成子群循环”不表示$\Q$本身循环。若$\Q=\langle q\rangle$且$q\ne0$，则$q/2$不是$q$的整数倍，矛盾。旧记号$\Q^+$在此指有理数加法群。


\subsection{习题8：\texorpdfstring{有限阶生成元与交换性推出有限性}{有限阶生成元与交换性推出有限性}}
\label{ex:2:8}

\paragraph{题目}
交换群$G$由有限多个有限阶元素生成，证明$G$有限。

\paragraph{考点与思路}
交换性允许集中同一生成元，有限阶允许把指数限制在有限范围内。

\paragraph{所用结论与原文依据}
\noteref{generated}；\noteref{powers}；\noteref{abelianpowers}。

% 原解来源：chapters/Part 2/chapter2-exercises.tex:128；第2章原\exref{2}{8}。

\begin{proof}
设$G=\langle a_1,\ldots,a_t\rangle$，记$n_i=\operatorname{ord}(a_i)<\infty$。由生成的定义，任意$g\in G$是有限多个$a_i$和$a_i^{-1}$的乘积。由于$G$交换，可重排这些因子，将同类因子合并，得到
\[
g=a_1^{k_1}\cdots a_t^{k_t},\qquad k_i\in\Z.
\]
对每个指数作带余除法$k_i=q_in_i+r_i$，其中$0\le r_i<n_i$。由$a_i^{n_i}=e$，有$a_i^{k_i}=a_i^{r_i}$。故映射
\[
\prod_{i=1}^t\{0,1,\ldots,n_i-1\}\longrightarrow G,
\qquad(r_1,\ldots,r_t)\longmapsto a_1^{r_1}\cdots a_t^{r_t}
\]
满射。定义域只有$\prod_i n_i$个元素，因而
\[
|G|\le\prod_{i=1}^t n_i<\infty.
\]
空生成集对应平凡群，也满足结论。
\end{proof}

\paragraph{复习提示}
上述表达式未必唯一，所以一般只能写上界。去掉交换性后，有限阶生成元之间交替相乘可能产生无限多个元素。


\subsection{习题9：\texorpdfstring{由幂集合刻画循环群}{由幂集合刻画循环群}}
\label{ex:2:9}

\paragraph{题目}
令$G^k=\{a^k:a\in G\}$，$k\geq1$。
$G$为有限循环群，当且仅当它的每个子群都等于某个$G^k$。

\paragraph{考点与思路}
先由幂集合的周期性证明群有限，再证明各西罗子群唯一且循环。

\paragraph{所用结论与原文依据}
\noteref{cyclic}；\noteref{powerorder}；\noteref{sylow1}；\noteref{sylow2}；\noteref{internal}；\noteref{coprime}。

% 原解来源：chapters/Part 2/chapter2-exercises.tex:148；第2章原\exref{2}{9}。

\begin{proof}
\textbf{必要性。}设$G=\langle a\rangle$为有限循环群，$H\leq G$。循环群的子群分类给出$H=\langle a^d\rangle$，其中$d\mid|G|$。又
\[
G^d=\{(a^t)^d:t\in\Z\}=\langle a^d\rangle=H.
\]
平凡子群对应$d=|G|$，所以每个子群均是某个$G^k$。

\textbf{充分性。}设每个子群均为某个$G^k$。对平凡子群，有正整数$E$使$G^E=\{e\}$。故$g^E=e$对每个$g\in G$成立，并有$G^{k+E}=G^k$。不同幂集合至多$E$个，因而不同子群也只有有限多个。每个循环子群的阶整除$E$，均有限；又
\[
G=\bigcup_{g\in G}\langle g\rangle,
\]
故$G$是有限多个有限循环子群的并，因而有限。

固定素数$p\mid|G|$。从全部$p$幂阶元素中选择阶最大的$a$，记$\ord(a)=p^r$。这种元素存在：\proofref{sylow1}{西罗第一定理}给出$p$阶子群，而$G$有限。题设给出$\langle a\rangle=G^k$，$k\geq1$。因为$a\in G^k$，可取$b\in G$使$b^k=a$。

若$p\mid k$，写$\ord(b)=p^su$、$p\nmid u$，以及$k=p^tv$、$p\nmid v$、$t\geq1$。\proofref{powerorder}{幂的阶公式}说明$b^k$的阶中$p$的指数为$\max(s-t,0)=r$，所以$s=r+t>r$。但$b^u$的阶恰为$p^s$，与$a$的选取矛盾。因此$p\nmid k$。

任取$p$幂阶元素$x$，设$\ord(x)=p^l$。存在整数$v$使$kv\equiv1\pmod{p^l}$，故
\[
x=(x^v)^k\in G^k=\langle a\rangle.
\]
所以所有$p$幂阶元素均属于$\langle a\rangle$。\proofref{sylow2}{西罗第二定理}保证有西罗$p$子群$P$包含$\langle a\rangle$；而$P$的每个元素都是$p$幂阶，反向也属于$\langle a\rangle$。故$P=\langle a\rangle$。同理任何西罗$p$子群都等于它，所以该西罗子群唯一、正规且循环。

对不同素数$p,q$，对应的正规西罗子群交平凡；两个因子元素的换位子同时属于两子群，故为$e$，所以两因子逐元素交换。它们的乘积是子群，反复使用第一章习题21的计数公式，其阶等于$|G|$，故等于$G$。各循环因子的阶两两互素，由原笔记互素阶交换元素乘积的阶的引理，各生成元的积具有$|G|$阶，因而生成$G$。所以$G$为有限循环群。

这里$k$限定为正整数。无限循环群的任何$G^k$都非平凡，不能表示平凡子群，所以结论应写作有限循环群版本。
\end{proof}

\paragraph{复习提示}
幂集合$G^k$一般未必为子群；本题仅由题设知道指定的幂集合等于某个子群。


\subsection{习题10：\texorpdfstring{两组对换生成$S_n$}{两组对换生成Sn}}
\label{ex:2:10}

\paragraph{题目}
证明$S_n$由$(12),(13),\ldots,(1n)$生成，也由$(12),(23),\ldots,(n-1\ n)$生成。

\paragraph{考点与思路}
每个置换都是对换的乘积；因此只需把任意对换表示成指定生成元的乘积。

\paragraph{所用结论与原文依据}
\noteref{threecycles}；\noteref{conjugation}。

% 原解来源：chapters/Part 2/chapter2-exercises.tex:200；第2章原\exref{2}{10}。

\begin{proof}
先考虑第一组。涉及点1的对换已经在列表中；若$i,j\ne1$且$i\ne j$，利用\proofref{conjugation}{共轭公式}
\[
u(ab)u^{-1}=(u(a)\ u(b))
\]
取$u=(1i)$、$(ab)=(1j)$，得
\[
(ij)=(1i)(1j)(1i).
\]
故第一组生成全部对换，进而生成$S_n$。

再考虑相邻对换。固定$i<j$。若$j=i+1$，$(ij)$已经是相邻对换。若$j>i+1$，令
\[
u=(i\ i+1)(i+1\ i+2)\cdots(j-2\ j-1).
\]
从右至左作用时，$j-1$逐步被送到$i$，而$j$不受影响，因此$u(j-1)=i$、$u(j)=j$。共轭公式给出
\[
(ij)=u(j-1\ j)u^{-1}.
\]
右边的$u$、$u^{-1}$和中间对换都由相邻对换组成。因此第二组也生成全部对换，从而生成$S_n$。
\end{proof}

\paragraph{复习提示}
必须说明“生成任意对换”为什么足够；共轭公式还将用于\exref{2}{11}、12。


\subsection{习题11：\texorpdfstring{两个生成元生成$S_n$}{两个生成元生成Sn}}
\label{ex:2:11}

\paragraph{题目}
当$n\ge2$时，证明$S_n=\langle(12),(12\cdots n)\rangle$。

\paragraph{考点与思路}
用长轮换的幂把$(12)$共轭成所有相邻对换。

\paragraph{所用结论与原文依据}
\noteref{conjugation}；\noteref{generated}；\exref{2}{10}。

% 原解来源：chapters/Part 2/chapter2-exercises.tex:221；第2章原\exref{2}{11}。

\begin{proof}
令$c=(12\cdots n)$，$K=\langle(12),c\rangle\le S_n$。对$1\le i<n$，有$c^{i-1}(1)=i$、$c^{i-1}(2)=i+1$，故
\[
c^{i-1}(12)c^{-(i-1)}=(i\ i+1)\in K.
\]
于是$K$包含\exref{2}{10}中生成$S_n$的全部相邻对换，从而$S_n\le K$。结合$K\le S_n$，即得$K=S_n$。$n=2$时两个指定生成元相同，仍生成$S_2$。
\end{proof}

\paragraph{复习提示}
生成集可以有重复元素；关键是其生成子群是否包含一组已知生成元。


\subsection{习题12：\texorpdfstring{两个生成元生成$A_n$}{两个生成元生成An}}
\label{ex:2:12}

\paragraph{题目}
当$n>2$为奇数时，证明
\[
A_n=\langle(123),(12\cdots n)\rangle;
\]
当$n>2$为偶数时，证明
\[
A_n=\langle(123),(23\cdots n)\rangle.
\]
\textbf{原题订正：}按教材书页98（PDF第108页）核对，原笔记将两个生成元抄为偶数时$(12\cdots n)^2$、奇数时$(34\cdots n)$；本题已恢复教材中的生成元。下面仍沿用原笔记的“共轭构造3-轮换，再逐次附加新点”的证明思路。

\paragraph{考点与思路}
先由共轭取得一列3-轮换，再说明这些轮换涉及的点彼此连接时，为什么能生成全部偶置换。

\paragraph{所用结论与原文依据}
\noteref{conjugation}；\noteref{threecycles}。

% 原解来源：chapters/Part 2/chapter2-exercises.tex:254；第2章原\exref{2}{12}。

\begin{proof}
先证明一个加入新点的结论。设$U$至少有三个点，某子群$K$已包含在$U$上作用、其余点固定的全部偶置换。

若$K$还含$(abz)$，其中$a,b\in U$、$z\notin U$，则$K$包含任意$(uvz)$，其中$u,v\in U$且$u\ne v$。理由如下：当$|U|\ge4$时，先取把有序点对$(a,b)$送至$(u,v)$的置换；如果它是奇置换，在目标中除$u,v$以外的两个点上复合一个对换即可把它改成偶置换。对$(abz)$作共轭就得到$(uvz)$。当$|U|=3$时，$U$上的3-轮换循环移动点对，可得到三个有序点对；再对$(abz)$取逆，可交换前两个位置，得到其余三个有序点对。因此此时结论也成立。

任何$U\cup\{z\}$上的3-轮换，要么只涉及$U$，要么涉及$z$和$U$中的两个点。以上已经把两类全部包含，所以$K$包含$U\cup\{z\}$上的全部偶置换。

若新增的3-轮换只与$U$交一个点，写为$(cde)$，其中$c\in U$、$d,e\notin U$。按照现有笔记的附加点引理，选$u\in U\setminus\{c\}$，用$U$内把$c$送到$u$的偶置换共轭，得到$(ude)\in K$。这样的偶置换存在：在$U$中另取不同于$c,u$的一点，选适当方向的3-轮换即可。

现在先把$\{c,d,e\}$当成已有点集：$(cde)$的幂已经给出该三点集上的全部偶置换，而$(ude)$含有其中的$d,e$两个点。由前面加入一个新点的结论，$K$包含$\{c,d,e,u\}$上的全部偶置换，特别包含$(cud)$与$(cue)$。

回到原点集$U$，轮换$(cud)$中$c,u$是旧点，故先把$d$加入，得到$U\cup\{d\}$上的全部偶置换；再用$(cue)$把$e$加入，得到$U\cup\{d,e\}$上的全部偶置换。因此，\textbf{从一个3-轮换起，逐次加入与已有点集有交的3-轮换，就能生成最终点集上的全部偶置换}。

\textbf{奇数$n$的情形。}令$c=(12\cdots n)$、$K=\langle(123),c\rangle$。轮换$c$长为奇数，符号为$(-1)^{n-1}=1$，故两个生成元均为偶置换，$K\le A_n$。对$j=0,1,\ldots,n-3$，\proofref{conjugation}{共轭公式}给出
\[
c^j(123)c^{-j}=(j+1\ j+2\ j+3)\in K.
\]
这些轮换依次为$(123),(234),\ldots,(n-2\ n-1\ n)$。从三点集$\{1,2,3\}$开始，下一轮换的前两个点已经出现，第三个点是新点，所以反复使用前面的附加一个新点结论，即得到全部$n$个点上的偶置换。因此$A_n\le K$，结合反向包含，得$K=A_n$。当$n=3$时，$c=(123)$，它本身已经生成$A_3$。

\textbf{偶数$n$的情形。}此时$n\ge4$。令$d=(23\cdots n)$、$K=\langle(123),d\rangle$。轮换$d$长为奇数$n-1$，符号为$(-1)^{n-2}=1$，故$K\le A_n$。

$d$固定1。对$j=0,1,\ldots,n-3$，有$d^j(2)=j+2$、$d^j(3)=j+3$，因此
\[
d^j(123)d^{-j}=(1\ j+2\ j+3)\in K.
\]
这些轮换依次为$(123),(134),(145),\ldots,(1\ n-1\ n)$。从$\{1,2,3\}$起，每次轮换中的1和另一个点属于已有点集，第三个点是新点。因此仍逐次用附加一个新点的结论，得到全部$n$个点上的偶置换，即$A_n\le K$。所以$K=A_n$。
\end{proof}

\paragraph{复习提示}
只说明生成元都是偶置换，只能得到生成子群包含于$A_n$。还须证明它足够大；这里的关键是将共轭构造与加入新点的论证结合。


\subsection{习题13：\texorpdfstring{长轮换的中心化子}{长轮换的中心化子}}
\label{ex:2:13}

\paragraph{题目}
若$\sigma=(12\cdots n)$，则$C_{S_n}(\sigma)=\langle\sigma\rangle$，其共轭类大小为$(n-1)!$。

\paragraph{考点与思路}
交换关系使置换完全由一个点的像决定，再用轨道稳定子定理计数。

\paragraph{所用结论与原文依据}
\noteref{centerizer}；\noteref{orbit}；\noteref{cycles}。

% 原解来源：chapters/Part 2/chapter2-exercises.tex:278；第2章原\exref{2}{13}。

\begin{proof}
若$\tau\in C_{S_n}(\sigma)$，则$\tau\sigma=\sigma\tau$。反复使用这一等式，对$j\geq0$有
\[
\tau(\sigma^j(1))=\sigma^j(\tau(1)).
\]
因为$\sigma$为$n$轮换，$1,\sigma(1),\ldots,\sigma^{n-1}(1)$穷尽全部$n$个点。因此$\tau$完全由$\tau(1)$决定。

存在唯一$0\leq k<n$使$\tau(1)=\sigma^k(1)$，代入上式得
\[
\tau(\sigma^j(1))=\sigma^{j+k}(1)=\sigma^k(\sigma^j(1)).
\]
所以$\tau=\sigma^k$。反之，$\sigma$的每个幂都与$\sigma$交换，故$C_{S_n}(\sigma)=\langle\sigma\rangle$，其阶为$n$。

对共轭作用，$\sigma$的稳定子就是其中心化子。由\proofref{orbit}{轨道稳定子定理}，$\sigma$的共轭类大小为
\[
[S_n:C_{S_n}(\sigma)]=\frac{n!}{n}=(n-1)!.
\]
\end{proof}

\paragraph{复习提示}
证明中心化子相等要给出两个包含；共轭类大小由中心化子的指数决定。


\subsection{习题14：\texorpdfstring{$S_5$的共轭类及正规子群}{S5的共轭类及正规子群}}
\label{ex:2:14}

\paragraph{题目}
列出$S_5$的所有共轭类，计算各类大小，并确定全部正规子群。

\paragraph{考点与思路}
按轮换类型列共轭类，再利用正规子群的阶及乘法封闭性筛选。

\paragraph{所用结论与原文依据}
\noteref{conjugation}；\noteref{cycles}；\noteref{normal}；\noteref{lagrange}；\noteref{sign}。

% 原解来源：chapters/Part 2/chapter2-exercises.tex:298；第2章原\exref{2}{14}。

\begin{proof}
共轭把轮换中的符号逐个替换，所以同一轮换类型的置换共轭，反之共轭也保持轮换类型。5的分拆对应七种类型，其代表与类大小为
\[
\begin{array}{c|rrrrrrr}
\text{代表}&e&(12)&(12)(34)&(123)&(123)(45)&(1234)&(12345)\\\hline
\text{类大小}&1&10&15&20&20&30&24
\end{array}
\]
计数依据如下：对换有$\binom52=10$个；双对换先选固定点，再把余下四点分成两对，有$5\cdot3=15$个；3-轮换先选三点，再选两种循环方向，有$\binom53\cdot2=20$个；$(3,2)$型同样选三点及方向，余下两点的对换唯一，亦有20个；4-轮换有$\binom54\cdot3!=30$个；5-轮换有$4!=24$个。总数为120，已穷尽$S_5$。

正规子群$N$必须含$e$，并为共轭类的并；其阶还须整除120。将类大小$10,15,20,20,30,24$分别记为$10\epsilon_1,15\epsilon_2,20\epsilon_3,20\epsilon_4,30\epsilon_5,24\epsilon_6$，其中$\epsilon_i\in\{0,1\}$。核对这$2^6$种和与120的因子，可能阶数只有$1,40,60,120$。其中40只能是$1+15+24$，60只能是$1+15+20+24$。

40阶候选含全部双对换，但
\[
((12)(34))((12)(35))=(354)
\]
是3-轮换，不在该候选内，所以不对乘法封闭。60阶候选同样含全部双对换，上式说明它还必须包含3-轮换类；因此两个大小为20的类中只能选择3-轮换类。所得集合恰为所有偶置换，即$A_5$，确为正规子群，因为它是符号同态的核。

故$S_5$的全部正规子群为$\{e\},A_5,S_5$。
\end{proof}

\paragraph{复习提示}
共轭类大小相加得到的阶符合整除条件，只能作为候选，还要检查乘法封闭。


\subsection{习题15：\texorpdfstring{商群中子群的交}{商群中子群的交}}
\label{ex:2:15}

\paragraph{题目}
设$N\trianglelefteq G$、$H_1,H_2\le G$，且$N\le H_1\cap H_2$。证明
\[
(H_1/N)\cap(H_2/N)\cong(H_1\cap H_2)/N.
\]
期中复习册题面要求同构；现有笔记通过陪集逐项检验证明两者在$G/N$中实际相等，以下保留此原有方法，从而得到所求同构。

\paragraph{考点与思路}
在商群$G/N$中逐个陪集检验。包含$N$的条件保证一个代表元在子群内时，整个陪集都在其中。

\paragraph{所用结论与原文依据}
\noteref{correspondence}；\noteref{normal}；\noteref{quotient}。

% 原解来源：chapters/Part 2/chapter2-exercises.tex:330；第2章原\exref{2}{15}。

\begin{proof}
因为$N\trianglelefteq G$且$N\le H_i$，$N$在每个$H_i$中也正规，各商群均有定义。对任意$g\in G$，证明
\[
gN\in H_i/N\iff g\in H_i.
\]
若$g\in H_i$，这是$H_i/N$的定义。反过来，若$gN=hN$，其中$h\in H_i$，则$h^{-1}g\in N\le H_i$，所以$g=h(h^{-1}g)\in H_i$。

因此对任意$gN\in G/N$，有
\begin{align*}
gN\in(H_1/N)\cap(H_2/N)
&\iff g\in H_1\text{ 且 }g\in H_2\\
&\iff g\in H_1\cap H_2\\
&\iff gN\in(H_1\cap H_2)/N.
\end{align*}
两边由同样的陪集组成，故相等。
\end{proof}

\paragraph{复习提示}
结论是在$G/N$中的实际相等，比“同构”更强。不能遗漏$N\le H_i$，否则$H_i/N$甚至未必有定义。


\subsection{习题16：\texorpdfstring{交子群的陪集}{交子群的陪集}}
\label{ex:2:16}

\paragraph{题目}
$H_1\cap H_2$的每个左陪集可写成$H_1$一个左陪集与$H_2$一个左陪集的交。

\paragraph{考点与思路}
先用同一个代表元写出所需交集，再说明非空交集的一般形式。

\paragraph{所用结论与原文依据}
\noteref{cosets}；\noteref{glaws}。

% 原解来源：chapters/Part 2/chapter2-exercises.tex:348；第2章原\exref{2}{16}。

\begin{proof}
任取$g\in G$。由集合元素的写法，
\[
g(H_1\cap H_2)=gH_1\cap gH_2.
\]
具体地，若$x=gh$且$h\in H_1\cap H_2$，则$x$同时属于$gH_1,gH_2$。反之，若$x=gh_1=gh_2$，其中$h_i\in H_i$，则由左消去律$h_1=h_2$，故$x\in g(H_1\cap H_2)$。这已证明题设结论。

若还考虑两个任意左陪集$aH_1,bH_2$，其交非空时取$x$属于交。由陪集相交即相等，有$aH_1=xH_1$、$bH_2=xH_2$，所以
\[
aH_1\cap bH_2=x(H_1\cap H_2).
\]
因此任意两个此类左陪集的交，或为空，或为交子群的一个左陪集。
\end{proof}

\paragraph{复习提示}
两个任意陪集的交可能为空，写成某个陪集前须检查非空。


\subsection{习题17：\texorpdfstring{有限指数子群的交}{有限指数子群的交}}
\label{ex:2:17}

\paragraph{题目}
若$[G:H_1],[G:H_2]$有限，则
$[G:H_1\cap H_2]\leq[G:H_1][G:H_2]$。

\paragraph{考点与思路}
把交子群的陪集映入两个陪集集合的笛卡尔积，并验证良定义和单射。

\paragraph{所用结论与原文依据}
\noteref{cosets}。

% 原解来源：chapters/Part 2/chapter2-exercises.tex:365；第2章原\exref{2}{17}。

\begin{proof}
记$\mathcal L_G(H)=\{gH:g\in G\}$为左陪集集合，不要求$H$正规。定义
\[
\theta:\mathcal L_G(H_1\cap H_2)\longrightarrow
\mathcal L_G(H_1)\times\mathcal L_G(H_2),\quad
g(H_1\cap H_2)\longmapsto(gH_1,gH_2).
\]
若$g(H_1\cap H_2)=g'(H_1\cap H_2)$，则$g^{-1}g'\in H_1\cap H_2$，故$gH_i=g'H_i$对$i=1,2$均成立。这说明映射不依赖代表元，因而良定义。

若$\theta(g(H_1\cap H_2))=\theta(g'(H_1\cap H_2))$，则$gH_i=g'H_i$，故$g^{-1}g'\in H_i$对$i=1,2$成立，从而$g^{-1}g'\in H_1\cap H_2$。所以源陪集相同，$\theta$为单射。

题设保证右侧有限，其元素个数为$[G:H_1][G:H_2]$。因此源陪集集合也有限，并有
\[
[G:H_1\cap H_2]\leq[G:H_1][G:H_2].
\]
\end{proof}

\paragraph{复习提示}
有限指数不意味着群本身有限；这里计数的是陪集集合。


\subsection{习题18：\texorpdfstring{有限指数给出正规子群或忠实表示}{有限指数给出正规子群或忠实表示}}
\label{ex:2:18}

\paragraph{题目}
有限群$G$有真子群$H$、$[G:H]=n>1$，则或有非平凡正规子群，其指数整除$n!$，或$G$同构于$S_n$的子群。

\paragraph{考点与思路}
沿原解构造左陪集作用，按作用的核是否平凡分类。

\paragraph{所用结论与原文依据}
\noteref{action}；\noteref{firstiso}；\noteref{lagrange}。

% 原解来源：chapters/Part 2/chapter2-exercises.tex:385；第2章原\exref{2}{18}。

\begin{proof}
令$G$作用于$H$的$n$个左陪集：$g\cdot(xH)=gxH$。该式良定义，且单位元及乘积的作用满足群作用公理，故给出同态$\rho:G\to S_n$。设$K=\ker\rho$，则$K\trianglelefteq G$。

若$k\in K$，则$k$特别固定陪集$H$，即$kH=H$，故$k\in H$。因此$K\leq H<G$，所以$K$一定为真子群。

若$K\ne\{e\}$，它是非平凡真正规子群。由\proofref{firstiso}{第一同构定理}及\proofref{lagrange}{拉格朗日定理}，
\[
[G:K]=|G/K|=|\rho(G)|\mid|S_n|=n!.
\]
若$K=\{e\}$，则$\rho$单射，所以$G\cong\rho(G)\leq S_n$。两种情况已穷尽，故结论成立。
\end{proof}

\paragraph{复习提示}
核不可能为整个群，因为它包含于真子群$H$。


\subsection{习题19：\texorpdfstring{最小素因子指数的子群正规}{最小素因子指数的子群正规}}
\label{ex:2:19}

\paragraph{题目}
若$p$为$|G|$的最小素因子，$[G:H]=p$，则$H\trianglelefteq G$。

\paragraph{考点与思路}
由陪集作用的像的阶，同时整除$|G|$和$p!$来限制像。

\paragraph{所用结论与原文依据}
\noteref{action}；\noteref{firstiso}；\noteref{lagrange}；\noteref{orbit}。

% 原解来源：chapters/Part 2/chapter2-exercises.tex:401；第2章原\exref{2}{19}。

\begin{proof}
令$G$作用于$H$的$p$个左陪集，记对应同态为$\rho:G\to S_p$，核为$K$、像为$Q$。与第二章习题18同样，$K\leq H$。\proofref{firstiso}{第一同构定理}及\proofref{lagrange}{拉格朗日定理}给出
\[
|Q|\mid|G|,\qquad |Q|\mid p!.
\]
因为$p$是$|G|$的最小素因子，$|Q|$的任何素因子都不小于$p$。另一方面$p!$中唯一不小于$p$的素因子是$p$，且它的指数为1。所以$|Q|$只能为1或$p$。

陪集作用传递：任意$xH,yH$，元素$yx^{-1}$把前者映到后者。故$Q$在这$p$个陪集上的轨道大小为$p$，由\proofref{orbit}{轨道稳定子定理}得$p\mid|Q|$。因此$|Q|=p$。

于是$[G:K]=p=[G:H]$。又$K\leq H$，指数乘法公式给出$[H:K]=1$，故$K=H$。因为$K$为核，$H\trianglelefteq G$。
\end{proof}

\paragraph{复习提示}
最小素因子的条件用来排除$p!$中的较小素因子。


\subsection{习题20：\texorpdfstring{$p^2$阶群的分类}{p2阶群的分类}}
\label{ex:2:20}

\paragraph{题目}
设$p$为素数，$|G|=p^2$。证明$G$交换，并证明$p^2$阶群恰有$C_{p^2}$与$C_p\times C_p$两种同构类型。

\paragraph{考点与思路}
先由中心非平凡证明交换，再按是否存在$p^2$阶元素分类。

\paragraph{所用结论与原文依据}
\noteref{pcenter}；\noteref{lagrange}；\noteref{cyclic}；\noteref{internal}。

% 原解来源：chapters/Part 2/chapter2-exercises.tex:421；第2章原\exref{2}{20}。

\begin{proof}
有限$p$群的中心非平凡，所以$|Z(G)|$只能为$p$或$p^2$。若为$p^2$，则$G=Z(G)$，已经交换。若为$p$，则$G/Z(G)$为$p$阶循环群。取生成元的提升$a\in G$，每个群元素都可写成$a^iz$，其中$z\in Z(G)$。因而任意两个元素满足
\begin{align*}
(a^iz)(a^jw)&=a^{i+j}zw\\
&=a^{j+i}wz=(a^jw)(a^iz).
\end{align*}
所以$G$仍交换。

若有$p^2$阶元素，则它生成$G$，故$G\cong C_{p^2}$。否则每个非单位元的阶都是$p$。选$a\ne e$，再选$b\notin\langle a\rangle$。两个$p$阶子群$\langle a\rangle,\langle b\rangle$的交只能为$\{e\}$：若交非平凡，任一非单位元便生成两个子群，使它们相等，与$b$的选取矛盾。

因为$G$交换，由第一章习题21，$|\langle a\rangle\langle b\rangle|=p^2$，故乘积等于$G$。映射$C_p\times C_p\to G$，$(\bar i,\bar j)\mapsto a^ib^j$为同态且满射；核平凡，因为$a^i=b^{-j}$必须属于两个因子的交。故$G\cong C_p\times C_p$。

这两个群因是否存在$p^2$阶元素而不同构，所以恰为两种类型。
\end{proof}

\paragraph{复习提示}
中心商群循环推出交换的计算在这里已经给出，不把后面的\exref{2}{39}作为未证前提。


\subsection{习题21：\texorpdfstring{非交换6阶群}{非交换6阶群}}
\label{ex:2:21}

\paragraph{题目}
每个非交换6阶群同构于$S_3$。

\paragraph{考点与思路}
唯一的3阶子群给出正规循环因子，再由2阶元素的共轭作用决定乘法。

\paragraph{所用结论与原文依据}
\noteref{sylow1}；\noteref{sylow3}；\noteref{cyclic}；\noteref{dihedral}；\noteref{iso}；\exref{1}{21}。

% 原解来源：chapters/Part 2/chapter2-exercises.tex:443；第2章原\exref{2}{21}。

\begin{proof}
\proofref{sylow3}{西罗第三定理}给出$n_3\mid2$、$n_3\equiv1\pmod3$，所以$n_3=1$。记唯一的3阶子群为$\langle a\rangle$，它正规。\proofref{sylow1}{西罗第一定理}给出2阶元素$b$。

两个子群的交平凡，由第一章习题21，它们的集合乘积有$3\cdot2=6$个元素，故为整个$G$。特别地$a,b$生成$G$。

共轭$b$保持$\langle a\rangle$，且保持$a$的阶，所以$bab^{-1}$只能为$a$或$a^{-1}$。若为$a$，则$a,b$交换，由它们生成的$G$也交换，与题设矛盾。因此
\[
a^3=b^2=e,\qquad bab^{-1}=a^{-1}.
\]
由此每个元素可唯一写成$a^ib^j$，$0\leq i<3$、$0\leq j<2$。唯一性来自两个不同陪集$\langle a\rangle$和$\langle a\rangle b$，各自有三个不同元素。

定义$a^ib^j\mapsto(123)^i(12)^j$。$S_3$中对应元素满足相同关系，所以化简两个正规形乘积所得公式一致，映射保乘法；两边均有六个不同正规形，故为双射。因此$G\cong S_3$。
\end{proof}

\paragraph{复习提示}
仅列出几个成立的关系不足以直接判定同构，需补充正规形和双射的说明。


\subsection{习题22：\texorpdfstring{15阶群}{15阶群}}
\label{ex:2:22}

\paragraph{题目}
全部15阶群同构于$C_{15}$。

\paragraph{考点与思路}
两种西罗子群都唯一，正规性和互素阶使其逐元素交换。

\paragraph{所用结论与原文依据}
\noteref{sylow3}；\noteref{normal}；\noteref{coprime}；\exref{1}{21}。

% 原解来源：chapters/Part 2/chapter2-exercises.tex:461；第2章原\exref{2}{22}。

\begin{proof}
由\proofref{sylow3}{西罗第三定理}，
\begin{align*}
n_5\mid3,\quad n_5\equiv1\pmod5&\quad\Longrightarrow\quad n_5=1,\\
n_3\mid5,\quad n_3\equiv1\pmod3&\quad\Longrightarrow\quad n_3=1.
\end{align*}
设这两个唯一的西罗子群为$P=\langle a\rangle$、$Q=\langle b\rangle$，其中$a$为3阶、$b$为5阶。它们均正规，且交平凡，因为交的阶同时整除3和5。

对$x\in P,y\in Q$，正规性说明$[x,y]$同时属于$P$和$Q$，所以$[x,y]=e$，即$xy=yx$。又第一章习题21给出$|PQ|=15$，故$G=PQ$。

因此$a,b$交换，且阶互素。由原笔记互素阶交换元素乘积的阶的引理，$\ord(ab)=15$，所以$ab$生成$G$，故$G\cong C_{15}$。
\end{proof}

\paragraph{复习提示}
两个循环子群的乘积未必循环；这里生成元的阶互素且彼此交换才推出循环。


\subsection{习题23：\texorpdfstring{10阶群}{10阶群}}
\label{ex:2:23}

\paragraph{题目}
证明10阶群恰有$C_{10}$与$D_5$两种同构类型。

\paragraph{考点与思路}
根据2阶元素在唯一5阶子群上的共轭作用分类。

\paragraph{所用结论与原文依据}
\noteref{sylow1}；\noteref{sylow3}；\noteref{coprime}；\noteref{dihedral}；\exref{1}{21}。

% 原解来源：chapters/Part 2/chapter2-exercises.tex:478；第2章原\exref{2}{23}。

\begin{proof}
由$n_5\mid2$及$n_5\equiv1\pmod5$，得$n_5=1$。设唯一的5阶子群为$\langle a\rangle$，取2阶元素$b$。两个子群交平凡，由计数可得$G=\langle a\rangle\langle b\rangle$，所以$a,b$生成$G$。

由于$\langle a\rangle$正规，$bab^{-1}=a^t$，其中$\gcd(t,5)=1$。由$b^2=e$，再次共轭得
\[
a=b(a^t)b^{-1}=a^{t^2},
\]
所以$t^2\equiv1\pmod5$。因5为素数，$5\mid(t-1)(t+1)$推出$t\equiv1$或$-1\pmod5$。

若$t\equiv1$，则$a,b$交换，互素阶乘积的阶公式给出$\ord(ab)=10$，故$G\cong C_{10}$。

若$t\equiv-1$，则$a^5=b^2=e$且$bab^{-1}=a^{-1}$。每个元素唯一写为$a^i$或$a^ib$，$0\leq i<5$。将它们分别映到$D_5$中的$T^i$和$T^iS$，由相同的关系核对乘法，得到同构$G\cong D_5$。

$C_{10}$交换，$D_5$不交换，所以两者不同构；$t$的两种可能已穷尽，分类完成。
\end{proof}

\paragraph{复习提示}
先证明共轭只剩两种，再分别识别两个群，才能说明分类没有遗漏。


\subsection{习题24：\texorpdfstring{不同素数乘积阶的群不单}{不同素数乘积阶的群不单}}
\label{ex:2:24}

\paragraph{题目}
若$|G|=pq$、$p<q$，则$G$不是单群。

\paragraph{考点与思路}
对较大素数的西罗子群计数，排除非唯一的可能。

\paragraph{所用结论与原文依据}
\noteref{sylow3}；\noteref{simple}。

% 原解来源：chapters/Part 2/chapter2-exercises.tex:496；第2章原\exref{2}{24}。

\begin{proof}
设$|G|=pq$，$p<q$为素数。西罗$q$子群的个数$n_q$满足
\[
n_q\mid p,\qquad n_q\equiv1\pmod q.
\]
所以$n_q$只能为1或$p$。若$n_q=p$，则$q\mid p-1$，但$0<p-1<q$，不可能。因此$n_q=1$，唯一的西罗$q$子群$Q$正规。

由于$1<q<pq$，$Q$既非平凡子群也非整个群，所以$G$不是单群。
\end{proof}

\paragraph{复习提示}
单性排除的是非平凡真正规子群，须交代候选子群既非1阶也非全群。


\subsection{习题25：\texorpdfstring{$p^2q$阶群有正规西罗子群}{p2q阶群有正规西罗子群}}
\label{ex:2:25}

\paragraph{题目}
设$p,q$为不同的素数，$|G|=p^2q$。证明$G$至少有一个正规的西罗子群。

\paragraph{考点与思路}
先用西罗数的整除和同余约束缩小到12阶，再直接数元素。

\paragraph{所用结论与原文依据}
\noteref{sylow3}；\noteref{lagrange}。

% 原解来源：chapters/Part 2/chapter2-exercises.tex:509；第2章原\exref{2}{25}。

\begin{proof}
若$n_p=1$，已得到正规西罗子群。否则$n_p\mid q$给出$n_p=q$，再由$n_p\equiv1\pmod p$得$q>p$。

若$n_q=1$，同样已完成。反设$n_q\ne1$。由$n_q\mid p^2$，它只能为$p$或$p^2$。因$p<q$，$p$不可能模$q$余1，故$n_q=p^2$，于是
\[
q\mid p^2-1=(p-1)(p+1).
\]
$q$为素数，必整除$p-1$或$p+1$；前者因$0<p-1<q$不可能。因此$q\mid p+1$。又$q>p$，只能$q=p+1$。两个相邻整数都为素数，只可能$p=2,q=3$，故剩下12阶情形。

此时反设给出$n_2=3$、$n_3=4$。不同3阶子群交平凡，所以四个3阶子群共有$4(3-1)=8$个3阶元素。群中余下恰有四个元素，包含$e$。

任一西罗2子群有四个元素，其中的元素阶整除4，不能是3阶元，故它只能恰好等于这四个剩余元素组成的集合。因此西罗2子群唯一，与$n_2=3$矛盾。故至少有一个西罗子群正规。
\end{proof}

\paragraph{复习提示}
12阶时要说明每个西罗2子群都被同一个四元素集合容纳，才能推出唯一。


\subsection{习题26：\texorpdfstring{两个非交换的有限例子}{两个非交换的有限例子}}
\label{ex:2:26}

\paragraph{题目}
分别构造所有元素满足$x^3=e$、及所有元素满足$x^4=e$的有限非交换群。

\paragraph{考点与思路}
保留原解的单位上三角矩阵群和二面体群两个构造。

\paragraph{所用结论与原文依据}
\noteref{subgroup}；\noteref{dihedral}。

% 原解来源：chapters/Part 2/chapter2-exercises.tex:530；第2章原\exref{2}{26}。

\begin{proof}
在$\F_3$上记
\[
X(a,b,c)=\begin{pmatrix}1&a&c\\0&1&b\\0&0&1\end{pmatrix},\qquad
H=\{X(a,b,c):a,b,c\in\F_3\}.
\]
直接计算可得
\begin{align*}
X(a,b,c)X(a',b',c')&=X(a+a',b+b',c+c'+ab'),\\
X(a,b,c)^{-1}&=X(-a,-b,ab-c).
\end{align*}
所以$H$含单位矩阵，对乘法和取逆封闭，因而是$\GL_3(\F_3)$的子群，阶为$3^3=27$。

每个$X\in H$可写为$E+N$，其中$N$严格上三角，故$N^3=0$。由于$E$与$N$交换，且域特征为3，
\[
X^3=(E+N)^3=E+3N+3N^2+N^3=E.
\]
然而
\begin{align*}
(E+E_{12})(E+E_{23})&=E+E_{12}+E_{23}+E_{13},\\
(E+E_{23})(E+E_{12})&=E+E_{12}+E_{23}.
\end{align*}
所以$H$非交换，是第一个例子。

第二个例子取$D_4$。由二面体群的元素形式，旋转是$T$的幂，阶整除4；反射$ST^k$满足$(ST^k)^2=e$。所以每个元素$x$均满足$x^4=e$。又原笔记已证明$D_4$阶为8且不交换，故满足要求。
\end{proof}

\paragraph{复习提示}
矩阵群的封闭性与逆元应先验证，再讨论其阶和元素的幂。


\subsection{习题27：\texorpdfstring{双陪集计数}{双陪集计数}}
\label{ex:2:27}

\paragraph{题目}
有限群中，$H,K\leq G,a\in G$，有
\[
|HaK|=|H|[K:a^{-1}Ha\cap K].
\]

\paragraph{考点与思路}
把双陪集分解成右陪集，再准确比较两个右陪集何时相等。

\paragraph{所用结论与原文依据}
\noteref{cosets}。

% 原解来源：chapters/Part 2/chapter2-exercises.tex:556；第2章原\exref{2}{27}。

\begin{proof}
令$L=a^{-1}Ha\cap K$。它是$K$的子群。双陪集可写为
\[
HaK=\bigcup_{k\in K}Hak,
\]
每个$Hak=H(ak)$都是$H$的一个右陪集。对$k_1,k_2\in K$，右陪集相等的判据给出
\begin{align*}
Hak_1=Hak_2
&\iff (ak_1)(ak_2)^{-1}\in H\\
&\iff ak_1k_2^{-1}a^{-1}\in H\\
&\iff k_1k_2^{-1}\in a^{-1}Ha\cap K=L\\
&\iff Lk_1=Lk_2.
\end{align*}
所以不同$Hak$的个数等于$L$在$K$中的右陪集数，即$[K:L]$。不同的$H$右陪集两两不交，每个均有$|H|$个元素。因此
\[
|HaK|=|H|[K:a^{-1}Ha\cap K].
\]
\end{proof}

\paragraph{复习提示}
这里使用右陪集的判据$Hx=Hy\iff xy^{-1}\in H$，次序不能写反。


\subsection{习题28：\texorpdfstring{西罗2子群只有2阶时的指数2子群}{西罗2子群只有2阶时的指数2子群}}
\label{ex:2:28}

\paragraph{题目}
设有限群$G$的西罗2子群阶为2。证明$G$存在指数2的子群。

\paragraph{考点与思路}
仍用左正则表示，把一个2阶元素对应为奇置换，再取符号同态的核。

\paragraph{所用结论与原文依据}
\noteref{cayley}；\noteref{sign}；\noteref{firstiso}。

% 原解来源：chapters/Part 2/chapter2-exercises.tex:577；第2章原\exref{2}{28}。

\begin{proof}
因为西罗2子群阶为2，可写$|G|=2m$，其中$m$为奇数。取该子群中的非单位元$a$，则$a^2=e$。原笔记的\proofref{cayley}{凯莱定理}给出左正则表示
\[
\rho:G\longrightarrow S(G),\qquad \rho(g)(x)=gx.
\]
若$ax=x$，右消去$x$得$a=e$，矛盾。因此$\rho(a)$没有不动点。又$a^2=e$，所以$\rho(a)$的每个轨道都有两个元素，正是$\{x,ax\}$。这些轨道把$G$分成$m$对，故$\rho(a)$为$m$个不相交对换的积。

由于$m$为奇数，$\sgn(\rho(a))=(-1)^m=-1$。因此复合同态$\varepsilon=\sgn\circ\rho:G\to\{1,-1\}$满射。令$N=\ker\varepsilon$，由\proofref{firstiso}{第一同构定理}，
\[
G/N\cong\{1,-1\},\qquad [G:N]=2.
\]
所以$G$有指数2的子群，且这个子群正规。
\end{proof}

\paragraph{复习提示}
构造符号同态后须证明它非平凡，才能确定核的指数恰为2。


\subsection{习题29：\texorpdfstring{大子集的乘积}{大子集的乘积}}
\label{ex:2:29}

\paragraph{题目}
设$G$为有限群，$A,B\subseteq G$，且$|A|+|B|>|G|$。证明$AB=G$，其中$AB=\{ab:a\in A,b\in B\}$。

\paragraph{考点与思路}
要把任意$g$写成$ab$，等价于让$A$和$gB^{-1}$有交点。

\paragraph{所用结论与原文依据}
\noteref{setproduct}；\noteref{glaws}。补充前置工具：有限集合并集的计数。

% 原解来源：chapters/Part 2/chapter2-exercises.tex:595；第2章原\exref{2}{29}。

\begin{proof}
固定$g\in G$。取逆映射和左乘$g$的映射都是双射，所以$|gB^{-1}|=|B|$。若$A\cap gB^{-1}=\varnothing$，它们的并集作为$G$的子集应满足
\[
|A|+|B|=|A\cup gB^{-1}|\le|G|,
\]
与题设矛盾。故存在$a\in A\cap gB^{-1}$。

由$a\in gB^{-1}$，可取$b\in B$使$a=gb^{-1}$，右乘$b$得$g=ab\in AB$。因为$g$任意，$G\subseteq AB$。又$A,B\subseteq G$且群运算封闭，$AB\subseteq G$，所以$AB=G$。
\end{proof}

\paragraph{复习提示}
严格大于不能改成大于等于。例如在二阶群$\{e,a\}$中取$A=B=\{e\}$，有$|A|+|B|=|G|$，但$AB\ne G$。


\subsection{习题30：\texorpdfstring{真子群的共轭不能覆盖有限群}{真子群的共轭不能覆盖有限群}}
\label{ex:2:30}

\paragraph{题目}
设$G$为有限群，$H<G$为真子群。证明
\[
G\ne\bigcup_{g\in G}gHg^{-1}.
\]

\paragraph{考点与思路}
共轭子群的数目至多为$[G:H]$；它们都包含单位元，因此并集的大小比直接相加至少少掉若干次重复。

\paragraph{所用结论与原文依据}
\noteref{orbit}；\noteref{normalizer}；\noteref{cosets}。补充前置工具：有限集合并集的上界。

% 原解来源：chapters/Part 2/chapter2-exercises.tex:610；第2章原\exref{2}{30}。

\begin{proof}
记$|G|=n$、$|H|=h$、$m=[G:H]=n/h>1$。$G$通过共轭作用在$H$的共轭子群集合上。稳定$H$的元素正是$N_G(H)$，所以不同共轭子群的数目为
\[
r=[G:N_G(H)].
\]
对每个$h_0\in H$，有$h_0Hh_0^{-1}=H$，故$H\le N_G(H)$，由指数公式得到$r\le m$。

把不同共轭子群列为$H_1,\ldots,H_r$。每个$H_i$有$h$个元素且含$e$，因此
\begin{align*}
\left|\bigcup_{g\in G}gHg^{-1}\right|
&=1+\left|\bigcup_{i=1}^r(H_i\setminus\{e\})\right|\\
&\le1+\sum_{i=1}^r|H_i\setminus\{e\}|\\
&=1+r(h-1)\\
&\le1+m(h-1)=n-m+1<n.
\end{align*}
并集元素数严格小于$|G|$，故不可能覆盖$G$。当$H=\{e\}$时，该估计仍成立。
\end{proof}

\paragraph{复习提示}
必须按“不同共轭子群”计数，不能把每个$g$都当成给出一个新子群。有限性用于基数估计。


\subsection{习题31：\texorpdfstring{$S_4$的西罗2子群及西罗3子群}{S4的西罗2子群及西罗3子群}}
\label{ex:2:31}

\paragraph{题目}
求$S_4$的全部西罗2子群与西罗3子群。

\paragraph{考点与思路}
8阶子群可由两两分块的稳定子构造；3阶子群由3-轮换生成，再用西罗数条件证明没有遗漏。

\paragraph{所用结论与原文依据}
\noteref{sylow1}；\noteref{sylow3}；\noteref{orbit}；\noteref{cycles}；\exref{1}{25}。

% 原解来源：chapters/Part 2/chapter2-exercises.tex:632；第2章原\exref{2}{31}。

\begin{proof}
因$|S_4|=24=2^3\cdot3$，西罗2子群的阶为8，西罗3子群的阶为3。

\textbf{西罗2子群。}将四个点分成两个无序二元组，有三种分块
\[
\mathcal P_1=\{\{1,2\},\{3,4\}\},\quad
\mathcal P_2=\{\{1,3\},\{2,4\}\},\quad
\mathcal P_3=\{\{1,4\},\{2,3\}\}.
\]
令$P_i=\{\sigma\in S_4:\sigma(\mathcal P_i)=\mathcal P_i\}$，即保持分块、允许交换两块的置换。稳定子在复合及取逆下封闭，故$P_i$是子群。可选择是否交换两块，再分别选择每块内部的对应，共有$2\cdot2\cdot2=8$种选择，所以$|P_i|=8$，均为西罗2子群。

三者彼此不同：$(12)\in P_1$而$(12)\notin P_2,P_3$；$(13)\in P_2$而$(13)\notin P_3$。又$n_2\mid3$，因此已有的三个子群正是全部。它们也可写为
\[
P_1=\langle(1324),(12)\rangle,\quad
P_2=\langle(1234),(13)\rangle,\quad
P_3=\langle(1243),(14)\rangle.
\]
以中间一组为例，$r=(1234)$、$s=(13)$满足$r^4=s^2=e$、$srs=r^{-1}$；八个$r^k,sr^k$互异，故它们生成整个$P_2$。其余两组按点的重新命名得到。因此三个子群都同构于正方形的8阶二面体群$D_4$。

\textbf{西罗3子群。}$S_4$的3阶元素只能是3-轮换。先选三个移动点，再选两种轮换方向，共有$\binom43\cdot2=8$个。一个3阶循环群恰有两个非单位元，且不同3阶子群交仅含$e$，因此恰有四个：
\[
\langle(123)\rangle,\quad\langle(124)\rangle,\quad
\langle(134)\rangle,\quad\langle(234)\rangle.
\]
\end{proof}

\paragraph{复习提示}
西罗子群的“阶”与“个数”分开计算；本题分别是阶8、个数3，以及阶3、个数4。


\subsection{习题32：\texorpdfstring{56阶与148阶群不单}{56阶与148阶群不单}}
\label{ex:2:32}

\paragraph{题目}
证明56阶群和148阶群都不是单群。

\paragraph{考点与思路}
148阶直接求出唯一37阶子群；56阶若7阶子群不唯一，就用元素计数迫使西罗2子群唯一。

\paragraph{所用结论与原文依据}
\noteref{sylow1}；\noteref{sylow3}；\noteref{lagrange}。

% 原解来源：chapters/Part 2/chapter2-exercises.tex:655；第2章原\exref{2}{32}。

\begin{proof}
\textbf{148阶。}设$|G|=148=4\cdot37$。西罗定理给出
\[
n_{37}\mid4,\qquad n_{37}\equiv1\pmod{37}.
\]
4的正因子只有1、2、4，只有1满足同余条件，所以$n_{37}=1$。唯一的37阶子群正规，且其阶既不为1也不为148，故$G$不单。

\textbf{56阶。}设$|G|=56=8\cdot7$。由$n_7\mid8$、$n_7\equiv1\pmod7$，得$n_7=1$或8。若$n_7=1$，唯一7阶子群就是非平凡真正规子群。

设$n_7=8$。不同7阶子群的交是各自的子群，阶整除7；若交非平凡，两者就相等。因此它们的非单位元互不重复，共有
\[
8(7-1)=48
\]
个7阶元素。全群除去这些元素后只剩8个元素，其中含$e$。

任取西罗2子群$P$，它有8个元素。由\proofref{lagrange}{拉格朗日定理}，$P$中元素阶只能为1、2、4、8，不能为7。因此$P$包含于上述剩余的8元集合。两者大小相同，必相等。每个西罗2子群都只能是该集合，所以$n_2=1$，$P$正规。又$1<8<56$，故$G$不单。
\end{proof}

\paragraph{复习提示}
不能在知道$n_7=8$后就断言不单；真正的正规子群是进一步计数得到的唯一西罗2子群。


\subsection{习题33：\texorpdfstring{弗拉蒂尼论证}{弗拉蒂尼论证}}
\label{ex:2:33}

\paragraph{题目}
有限群$G$中，$N\trianglelefteq G$，$P$是$N$的西罗$p$子群，则$G=N N_G(P)$。

\paragraph{考点与思路}
在正规子群$N$内使用西罗共轭定理，将任意$g$分解为两个指定因子的乘积。

\paragraph{所用结论与原文依据}
\noteref{sylow2}；\noteref{normalizer}。

% 原解来源：chapters/Part 2/chapter2-exercises.tex:677；第2章原\exref{2}{33}。

\begin{proof}
任取$g\in G$。由于$N\trianglelefteq G$，$gPg^{-1}\leq N$；共轭保持阶，故$gPg^{-1}$与$P$同阶，也是$N$的西罗$p$子群。

在群$N$内使用\proofref{sylow2}{西罗第二定理}，存在$n\in N$使
\[
gPg^{-1}=nPn^{-1}.
\]
于是$(n^{-1}g)P(n^{-1}g)^{-1}=P$。按正规化子的定义，$n^{-1}g\in N_G(P)$，从而$g=n(n^{-1}g)\in NN_G(P)$。由于$g$任意，$G\subseteq NN_G(P)$；反过来，两个因子均在$G$中，故$NN_G(P)\subseteq G$。因此$G=NN_G(P)$。
\end{proof}

\paragraph{复习提示}
共轭元$n$必须取在$N$中，故西罗共轭定理是在群$N$内使用。


\subsection{习题34：\texorpdfstring{正规子群与商群的西罗子群}{正规子群与商群的西罗子群}}
\label{ex:2:34}

\paragraph{题目}
设$G$为有限群，$P$为其西罗$p$子群，$N\trianglelefteq G$。证明：
\begin{enumerate}
\item $P\cap N$为$N$的西罗$p$子群；
\item $PN/N$为$G/N$的西罗$p$子群。
\end{enumerate}

\paragraph{考点与思路}
先用共轭使$N$的一个西罗子群落入$P$，再计算商群中的阶。

\paragraph{所用结论与原文依据}
\noteref{sylow2}；\noteref{sylowdef}；\noteref{secondiso}。

% 原解来源：chapters/Part 2/chapter2-exercises.tex:695；第2章原\exref{2}{34}。

\begin{proof}
\begin{enumerate}
\item 取$N$的西罗$p$子群$Q$。把$Q$视为$G$的$p$子群，由\proofref{sylow2}{西罗第二定理}可知它包含于$G$的某个西罗$p$子群$P_1$。同一定理又给出$g\in G$使$gP_1g^{-1}=P$。

因为$N$正规，$gQg^{-1}\leq N$；又$Q\leq P_1$，故$gQg^{-1}\leq P$。所以
\[
gQg^{-1}\leq P\cap N.
\]
左边与$Q$同阶，已经是$N$的西罗$p$子群。右边是$N$的$p$子群，阶不能更大。因此两者阶相同，$P\cap N$也是$N$的西罗$p$子群。

\item 写$|G|=p^rm$、$|N|=p^sn$，其中$p\nmid mn$。由第一项，$|P\cap N|=p^s$。因$N$正规，$PN$为子群，商群$PN/N$有定义；\proofref{secondiso}{第二同构定理}给出
\[
PN/N\cong P/(P\cap N),\qquad |PN/N|=p^{r-s}.
\]
另一方面，$|G/N|=|G|/|N|$中素因子$p$的指数恰为$r-s$。所以$PN/N$具有商群阶中的全部$p$因子，按定义是$G/N$的西罗$p$子群。
\end{enumerate}
\end{proof}

\paragraph{复习提示}
第一项不要求$P\cap N$正规；它的西罗性质通过阶数达到最大值证明。


\subsection{习题35：\texorpdfstring{子群的西罗子群可由共轭后相交得到}{子群的西罗子群可由共轭后相交得到}}
\label{ex:2:35}

\paragraph{题目}
若$H\leq G$，$P$是$G$的西罗$p$子群，则存在$a\in G$使$aPa^{-1}\cap H$是$H$的西罗$p$子群。

\paragraph{考点与思路}
先使$H$的西罗子群包含于$P$的某个共轭，再比较交子群的阶。

\paragraph{所用结论与原文依据}
\noteref{sylow1}；\noteref{sylow2}；\noteref{lagrange}。

% 原解来源：chapters/Part 2/chapter2-exercises.tex:714；第2章原\exref{2}{35}。

\begin{proof}
取$H$的西罗$p$子群$Q$。$Q$也是$G$的$p$子群，由\proofref{sylow2}{西罗第二定理}，存在$a\in G$使$Q\leq aPa^{-1}$。因此
\[
Q\leq aPa^{-1}\cap H\leq H.
\]
$aPa^{-1}\cap H$作为$p$群$aPa^{-1}$的子群，其阶为$p$的幂；又它属于$H$，所以其阶至多为$|H|$中最大的$p$幂，即$|Q|$。结合它包含$Q$，得两者阶相同，故
\[
aPa^{-1}\cap H=Q.
\]
所以该交为$H$的西罗$p$子群。
\end{proof}

\paragraph{复习提示}
不必要求$H$正规；共轭$P$的位置可以随$H$选择。


\subsection{习题36：\texorpdfstring{$\GL_n(\F_p)$的西罗$p$子群}{GLn(Fp)的西罗p子群}}
\label{ex:2:36}

\paragraph{题目}
设$p$为素数、$n\ge1$，求$\GL_n(\F_p)$的一个西罗$p$子群，并确定其正规化子及西罗$p$子群的个数。

\paragraph{考点与思路}
先构造单位上三角群并计算其阶，再用有序基计数核对它恰好具有群阶中的全部$p$因子；通过公共不动向量逐层恢复标准子空间链，确定正规化子。

\paragraph{所用结论与原文依据}
\noteref{sylow1}；\noteref{sylow2}；\noteref{sylow3}；\noteref{orbit}；\exref{1}{44}、\exref{1}{38}。

% 原解来源：chapters/Part 2/chapter2-exercises.tex:732；第2章原\exref{2}{36}。

\begin{proof}
记$G=\GL_n(\F_p)$，按照现有笔记先取
\[
U=\{(a_{ij}):a_{ii}=1,\ a_{ij}=0\ (i>j)\}.
\]
两个这样的矩阵之积仍为上三角且对角线全1。若$u=I+T\in U$，则$T$严格上三角，故$T^n=0$，并且
\[
u^{-1}=I-T+T^2-\cdots+(-1)^{n-1}T^{n-1}\in U.
\]
因而$U$是子群。严格上三角位置有$n(n-1)/2$个，每个位置可独立取$p$个值，故$|U|=p^{n(n-1)/2}$。

接着展开现有笔记所用的可逆矩阵计数公式。可逆矩阵的列组成有序基：第一列有$p^n-1$种选择；已选$k$个线性无关列时，它们张成的空间有$p^k$个向量，下一列有$p^n-p^k$种选择。因此
\begin{align*}
|G|&=\prod_{k=0}^{n-1}(p^n-p^k)\\
&=p^{0+1+\cdots+(n-1)}\prod_{i=1}^n(p^i-1)\\
&=p^{n(n-1)/2}\prod_{i=1}^n(p^i-1).
\end{align*}
每个$p^i-1$均不被$p$整除，所以$|G|$中最大的$p$幂因子是$p^{n(n-1)/2}$，恰好等于$|U|$，故$U$为西罗$p$子群。

令$B$为全体可逆上三角矩阵组成的群。上三角矩阵之积及逆仍为上三角，对角线在乘法下逐项相乘。因此$bub^{-1}$的对角线全1，得到$B\le N_G(U)$。

反过来，设$g\in N_G(U)$。当$n\ge2$时，以$e_i$表示第$i$个标准基向量，$E_{ij}$表示仅第$(i,j)$分量为1的矩阵。对$i<j$有$I+E_{ij}\in U$，且
\[
(I+E_{ij})v=v+v_je_i.
\]
若$v$被$U$的每个元素固定，取$i=1$、$j=2,\ldots,n$，就得$v_2=\cdots=v_n=0$。反过来$e_1$的所有倍数确实被全部$U$固定。因此公共不动空间为$L_1=\langle e_1\rangle$。

正规化条件保证$g$保持这个空间：若$v\in L_1$、$u\in U$，则$g^{-1}ug\in U$，所以
\[
u(gv)=g(g^{-1}ug)v=gv.
\]
于是$g(L_1)=L_1$。在商空间$\F_p^n/L_1$上，$U$的像恰为关于基$e_2+L_1,\ldots,e_n+L_1$的单位上三角群，而$g$诱导的映射仍正规化这个像。重复上述公共不动空间论证，或对维数归纳，可得$g$保持整个链
\[
L_1\subset L_2\subset\cdots\subset L_n,
\qquad L_i=\langle e_1,\ldots,e_i\rangle.
\]
这意味着第$i$列$g(e_i)$只能出现在前$i$个坐标位置，因此$g$上三角，即$g\in B$。所以$N_G(U)=B$。

$B$的每个对角线分量有$p-1$种选择，严格上三角分量仍有$p^{n(n-1)/2}$种选择，故
\[
|B|=(p-1)^np^{n(n-1)/2}.
\]
由共轭作用和\proofref{sylow2}{西罗共轭定理}，全部西罗$p$子群数为
\[
n_p=[G:N_G(U)]
=\frac{\prod_{i=1}^n(p^i-1)}{(p-1)^n}.
\]
当$n=1$时$U=\{1\}$、$B=G$，正规化子与计数公式仍成立。
\end{proof}

\paragraph{复习提示}
构造了具有正确阶的子群才算证明存在一个西罗子群。求个数时应除以正规化子阶，而非直接除以西罗子群阶。


\subsection{习题37：\texorpdfstring{非交换$p^3$阶群的换位子群}{非交换p3阶群的换位子群}}
\label{ex:2:37}

\paragraph{题目}
若$|G|=p^3$且$G$非交换，则$G^{(1)}=Z(G)$，两者阶为$p$。

\paragraph{考点与思路}
先确定中心的阶，再利用交换商群的判别限制换位子群。

\paragraph{所用结论与原文依据}
\noteref{pcenter}；\noteref{abelianization}；\noteref{lagrange}；\exref{2}{20}。

% 原解来源：chapters/Part 2/chapter2-exercises.tex:788；第2章原\exref{2}{37}。

\begin{proof}
因为$G$为非平凡$p$群，$Z(G)\ne\{e\}$。由\proofref{lagrange}{拉格朗日定理}，其阶只能为$p,p^2,p^3$。

若$|Z(G)|=p^3$，则$G$交换，与题设矛盾。若$|Z(G)|=p^2$，则$G/Z(G)$为$p$阶循环群。取其生成元的提升$a$，任意两个群元素可写成$a^iz,a^jw$，其中$z,w\in Z(G)$，并有
\[
(a^iz)(a^jw)=a^{i+j}zw=a^{j+i}wz=(a^jw)(a^iz).
\]
仍推出$G$交换，矛盾。因此$|Z(G)|=p$。

商群$G/Z(G)$阶为$p^2$，由第二章习题20，它为交换群。将原笔记的交换商群判别用于自然满同态$G\to G/Z(G)$，得$G^{(1)}\leq Z(G)$。由于$G$非交换，$G^{(1)}\ne\{e\}$；而$p$阶循环群$Z(G)$没有非平凡真子群，所以$G^{(1)}=Z(G)$，两者阶均为$p$。
\end{proof}

\paragraph{复习提示}
先在本题补写中心商循环的交换性计算，避免依赖尚未出现的\exref{2}{39}。


\subsection{习题38：\texorpdfstring{$p$群中的$p$阶正规子群在中心内}{p群中的p阶正规子群在中心内}}
\label{ex:2:38}

\paragraph{题目}
设$G$为有限$p$群，$N\trianglelefteq G$且$|N|=p$，证明$N\le Z(G)$。

\paragraph{考点与思路}
共轭给出从$p$群到$p-1$阶自同构群的同态，其像的阶只能为1。

\paragraph{所用结论与原文依据}
\noteref{cyclicaut}；\noteref{firstiso}；\noteref{lagrange}。

% 原解来源：chapters/Part 2/chapter2-exercises.tex:805；第2章原\exref{2}{38}。

\begin{proof}
因$N$正规，每个$g\in G$都通过共轭给出$N$的自同构。定义
\[
\rho:G\longrightarrow\Aut(N),\qquad
\rho(g)(x)=gxg^{-1}\quad(x\in N).
\]
共轭映射可逆且保持乘法；又$c_gc_h=c_{gh}$，故$\rho$为同态。

$N$阶为素数，任一非单位元$a$均生成$N$。自同构由$a$的像唯一决定；像可以是$a,a^2,\ldots,a^{p-1}$中任意一个，每个选择都给出自同构。因此$|\Aut(N)|=p-1$。

由\proofref{firstiso}{第一同构定理}及\proofref{lagrange}{拉格朗日定理}，$|\rho(G)|$既整除$|G|=p^r$，又整除$p-1$。这两个整数互素，所以$|\rho(G)|=1$。故$gxg^{-1}=x$对所有$g\in G,x\in N$成立，也就是$xg=gx$。因此$N\le Z(G)$。
\end{proof}

\paragraph{复习提示}
正规性本身只保证共轭后仍在$N$；这里还利用阶数互素，才把“仍在$N$”加强为“每个元素都不变”。


\subsection{习题39：\texorpdfstring{中心商群循环则原群交换}{中心商群循环则原群交换}}
\label{ex:2:39}

\paragraph{题目}
若$G/Z(G)$是循环群，证明$G$为交换群。

\paragraph{考点与思路}
同一个陪集生成元的幂提供非中心部分，其余因子都可自由交换。

\paragraph{所用结论与原文依据}
\noteref{center}；\noteref{cosets}；\noteref{powers}。

% 原解来源：chapters/Part 2/chapter2-exercises.tex:822；第2章原\exref{2}{39}。

\begin{proof}
设$G/Z(G)=\langle aZ(G)\rangle$。任意$x\in G$的陪集都有形式$xZ(G)=a^mZ(G)$，其中$m\in\Z$，故$a^{-m}x\in Z(G)$。因此任意$x,y\in G$可写成
\[
x=a^mz,\qquad y=a^nw,
\qquad z,w\in Z(G),\ m,n\in\Z.
\]
中心元与所有元素交换，故
\begin{align*}
xy&=a^mza^nw\\
&=a^ma^nzw\\
&=a^{m+n}zw\\
&=a^{n+m}wz\\
&=a^nwa^mz=yx.
\end{align*}
$x,y$任意，所以$G$交换。上述论证不要求$G$有限，也包括商群平凡的情况。
\end{proof}

\paragraph{复习提示}
结论实际上进一步给出$Z(G)=G$，所以中心商群若循环，就只能是平凡群。


\subsection{习题40：\texorpdfstring{有限群嵌入某个交错群}{有限群嵌入某个交错群}}
\label{ex:2:40}

\paragraph{题目}
证明任意有限群都同构于某个交错群的子群。

\paragraph{考点与思路}
先用左乘作用嵌入对称群，再添两个点，用一个额外对换补偿奇偶性。

\paragraph{所用结论与原文依据}
\noteref{cayley}；\noteref{sign}；\noteref{firstiso}。

% 原解来源：chapters/Part 2/chapter2-exercises.tex:840；第2章原\exref{2}{40}。

\begin{proof}
设$|G|=m$。对每个$g\in G$，定义左乘置换$L_g:G\to G$、$x\mapsto gx$。$L_{g^{-1}}$为其逆，且$L_gL_h=L_{gh}$，所以$g\mapsto L_g$是到$m$点对称群的同态。若$L_g$为恒等映射，则$g=L_g(e)=e$，因此它单射。给$G$的元素编号后得到$G\hookrightarrow S_m$。

把$S_m$中置换延拓到$m+2$个点，令其固定$m+1,m+2$。令$t=(m+1\ m+2)$。对$\sigma\in S_m$，设$\varepsilon(\sigma)\in\{0,1\}$表示奇偶性，偶时为0，奇时为1，定义
\[
\iota(\sigma)=\sigma t^{\varepsilon(\sigma)}.
\]
因为$t$是奇置换，$\iota(\sigma)$总是偶置换，即属于$A_{m+2}$。

$t$只作用于最后两个点，因而与$S_m$中的元素交换。又奇偶性在复合下模2相加，所以
\begin{align*}
\iota(\sigma)\iota(\tau)
&=\sigma t^{\varepsilon(\sigma)}\tau t^{\varepsilon(\tau)}\\
&=\sigma\tau t^{\varepsilon(\sigma)+\varepsilon(\tau)}\\
&=\sigma\tau t^{\varepsilon(\sigma\tau)}
=\iota(\sigma\tau).
\end{align*}
若$\iota(\sigma)=e$，将其限制于前$m$个点，就有$\sigma=e$，故$\iota$单射。复合$G\hookrightarrow S_m\hookrightarrow A_{m+2}$即得所求。
\end{proof}

\paragraph{复习提示}
不能直接说$S_m\le A_{m+2}$；必须指定这个保持乘法的嵌入。额外对换恰好抵消原置换的奇偶性。


\subsection{习题41：\texorpdfstring{正规子群与换位子群交平凡时在中心内}{正规子群与换位子群交平凡时在中心内}}
\label{ex:2:41}

\paragraph{题目}
若$N\trianglelefteq G$且$N\cap G^{(1)}=\{e\}$，证明$N\le Z(G)$。

\paragraph{考点与思路}
任取一个$N$中元素与任意群元的换位子，证明它同时落在$N$和导群中。

\paragraph{所用结论与原文依据}
\noteref{normal}；\noteref{derived}。

% 原解来源：chapters/Part 2/chapter2-exercises.tex:860；第2章原\exref{2}{41}。

\begin{proof}
取$n\in N$、$g\in G$。正规性给出$g^{-1}ng\in N$；因为$n^{-1}\in N$且$N$对乘法封闭，
\[
[n,g]=n^{-1}g^{-1}ng=n^{-1}(g^{-1}ng)\in N.
\]
另一方面，$G^{(1)}$由全部换位子生成，故$[n,g]\in G^{(1)}$。于是$[n,g]\in N\cap G^{(1)}=\{e\}$，即
\[
n^{-1}g^{-1}ng=e.
\]
左乘$n$再左乘$g$得到$ng=gn$。因为$g$任意，$n\in Z(G)$；因为$n$任意，$N\le Z(G)$。
\end{proof}

\paragraph{复习提示}
正规性用于证明换位子属于$N$，交平凡用于证明换位子等于$e$，两个条件分别承担不同步骤。


\subsection{习题42：\texorpdfstring{可解群的子群与商群仍可解}{可解群的子群与商群仍可解}}
\label{ex:2:42}

\paragraph{题目}
证明可解群的子群与商群仍然可解。

\paragraph{考点与思路}
子群的导出列包含在原群的导出列中；满同态则把导出列逐项映满到像群的导出列。

\paragraph{所用结论与原文依据}
\noteref{derivedseries}；\noteref{derived}；\noteref{solvabledef}。

% 原解来源：chapters/Part 2/chapter2-exercises.tex:873；第2章原\exref{2}{42}。

\begin{proof}
设$G^{(r)}=\{e\}$。

\textbf{子群。}若$H\le G$，则$H$中任意两个元素的换位子也是$G$中的换位子，因此$H^{(1)}\le G^{(1)}$。逐次使用同样的理由，归纳得到$H^{(i)}\le G^{(i)}$对所有$i\ge0$成立。特别地，
\[
H^{(r)}\le G^{(r)}=\{e\},
\]
所以$H$可解。

\textbf{商群。}设$N\trianglelefteq G$，令$\pi:G\twoheadrightarrow G/N$为自然满同态。任意同态都满足
\[
\pi([x,y])=[\pi(x),\pi(y)].
\]
这给出$\pi(G^{(1)})\le(G/N)^{(1)}$。反过来，商群任意两个元素都能取原像，故商群的每个换位子都是原群一个换位子的像；这些换位子生成商群导群，所以反向包含也成立。于是$\pi(G^{(1)})=(G/N)^{(1)}$。

把此结论反复应用于$\pi$在各导群上的限制，得到
\[
(G/N)^{(i)}=\pi(G^{(i)}).
\]
取$i=r$，右边为$\{N\}$，故$G/N$可解。
\end{proof}

\paragraph{复习提示}
子群情形一般是包含关系；商群情形因映射满而得到相等关系。不能把“导群元素”都误认为单个换位子，它们还包括换位子的有限乘积。


\subsection{习题43：\texorpdfstring{$p^2q$阶群可解}{p2q阶群可解}}
\label{ex:2:43}

\paragraph{题目}
设$p,q$为不同素数，$|G|=p^2q$，证明$G$可解。

\paragraph{考点与思路}
沿用现有笔记的顺序：由\exref{2}{25}得到正规西罗子群，由\exref{2}{20}判断它与商群均交换，最后应用可解群的扩张性质。下面将所依赖的两道习题及扩张性质一并展开。

\paragraph{所用结论与原文依据}
\noteref{solvable}；\noteref{sylow3}；\noteref{pcenter}；\exref{2}{25}、\exref{2}{20}。

% 原解来源：chapters/Part 2/chapter2-exercises.tex:891；第2章原\exref{2}{43}。

\begin{proof}
\textbf{第一步：至少有一个西罗子群正规。}设$n_p,n_q$分别为西罗子群数。若$n_p=1$或$n_q=1$，已经得到结论。反设两者都大于1。由$n_p\mid q$，必有$n_p=q$；由$n_p\equiv1\pmod p$，得$q>p$。

又$n_q\mid p^2$且$n_q>1$，所以$n_q=p$或$p^2$。因为$p<q$，整数$p$不可能模$q$余1，故$n_q=p^2$。于是
\[
q\mid p^2-1=(p-1)(p+1).
\]
$q$为素数，故整除其中一个因子。$q>p$排除$q\mid p-1$，所以$q\mid p+1$。由于$p<q\le p+1$，有$q=p+1$。相邻两个整数都是素数，只可能$p=2,q=3$。

因此只剩$|G|=12$且$n_2=3,n_3=4$的情形。四个3阶子群互交仅含$e$，共含$4\cdot2=8$个3阶元素。其余元素只有四个，含$e$。每个西罗2子群有四个元素，且不能含3阶元素，所以必等于这四元集合，因而西罗2子群唯一，与$n_2=3$矛盾。第一步得证。

\textbf{第二步：阶为$p^2$的群交换。}先回顾$p$群的类方程：非中心共轭类的大小是大于1的$p$幂，因此
\[
|P|=|Z(P)|+\sum\text{非中心共轭类大小}
\]
说明非平凡有限$p$群的中心阶被$p$整除，中心非平凡。若$|P|=p^2$，则$|Z(P)|=p$或$p^2$。若为$p$，商群$P/Z(P)$阶为$p$，故循环；由\exref{2}{39}，$P$交换，反而有$Z(P)=P$，与中心阶为$p$矛盾。所以必有$Z(P)=P$，$P$交换。

\textbf{第三步：应用可解群的扩张性质。}取第一步给出的正规西罗子群$N$。若$|N|=p^2$，第二步说明$N$交换，而$G/N$为$q$阶循环群；若$|N|=q$，则$N$循环，而$G/N$阶为$p^2$，同样交换。因此两种情况下$N$与$G/N$均可解。

现有笔记“可解群的子群、同态像及扩张”命题说明：正规子群与相应商群均可解时，原群可解。在这里使用该命题，即得$G$可解。为补全使用理由，按该命题的原证明，设$N^{(s)}=\{e\}$、$(G/N)^{(r)}=\{N\}$。对自然满同态$\pi:G\to G/N$，\exref{2}{42}的导群像公式给出
\[
\pi(G^{(r)})=(G/N)^{(r)}=\{N\},
\]
所以$G^{(r)}\le N$。再逐次取$s$次导群，得到
\[
G^{(r+s)}=\bigl(G^{(r)}\bigr)^{(s)}\le N^{(s)}=\{e\}.
\]
这正是扩张性质推出$G$可解的过程。本题中$N$和$G/N$均交换，可取$r=s=1$。
\end{proof}

\paragraph{复习提示}
本题并未要求$p<q$；证明里只有在反设两个西罗子群都不正规时，才推出$q>p$。应特别记住12阶的计数补充。


\subsection{习题44：\texorpdfstring{有限$p$群可解}{有限p群可解}}
\label{ex:2:44}

\paragraph{题目}
证明任意有限$p$群都是可解群。

\paragraph{考点与思路}
沿用现有笔记的中心商群归纳：类方程给出非平凡中心，商去中心后阶数下降；中心与中心商群均可解，再应用可解群的扩张性质。

\paragraph{所用结论与原文依据}
\noteref{classeq}；\noteref{pcenter}；\noteref{solvable}。

% 原解来源：chapters/Part 2/chapter2-exercises.tex:903；第2章原\exref{2}{44}。

\begin{proof}
对群阶归纳。平凡群交换，可解。设结论对阶小于$|G|$的$p$群成立，且$G$非平凡。

由\exref{2}{43}中类方程的论证，$Z(G)\ne\{e\}$。若$Z(G)=G$，则$G$交换，$G^{(1)}=\{e\}$，结论成立。

以下设$Z(G)$为真子群。因$|Z(G)|>1$，商群$G/Z(G)$仍是$p$群，且阶严格小于$|G|$。由归纳假设，$G/Z(G)$可解。另一方面，$Z(G)$交换，故也可解；中心又是$G$的正规子群。因此对$N=Z(G)$应用可解群的扩张性质，就得到$G$可解，归纳完成。

这里的扩张性质正是\exref{2}{43}第三步按现有笔记命题所展开的结论：若$(G/Z(G))^{(r)}$平凡，则$G^{(r)}\le Z(G)$；由于$Z(G)^{(1)}=\{e\}$，再取一次导群即有$G^{(r+1)}=\{e\}$。这补明了原文“由正规子群与商群可解推出原群可解”的理由。
\end{proof}

\paragraph{复习提示}
$p$群不一定交换，但中心一定非平凡。归纳所需的是商群阶数变小，不能仅说“商群仍是$p$群”就直接套结论。


\subsection{习题45：\texorpdfstring{两个可解商群给出交核的可解商群}{两个可解商群给出交核的可解商群}}
\label{ex:2:45}

\paragraph{题目}
若$H,K\trianglelefteq G$且$G/H,G/K$均可解，证明$G/(H\cap K)$可解。

\paragraph{考点与思路}
将两种商映射合并成到直积的映射，其核正好是交。再证明两个可解群的直积可解。

\paragraph{所用结论与原文依据}
\noteref{firstiso}；\noteref{derivedseries}；\noteref{solvable}；\noteref{internal}；\exref{2}{42}。

% 原解来源：chapters/Part 2/chapter2-exercises.tex:918；第2章原\exref{2}{45}。

\begin{proof}
定义
\[
\varphi:G\longrightarrow(G/H)\times(G/K),
\qquad g\longmapsto(gH,gK).
\]
两坐标都是同态，所以$\varphi$为同态。且
\[
g\in\ker\varphi\iff gH=H\text{ 且 }gK=K
\iff g\in H\cap K.
\]
\proofref{firstiso}{第一同构定理}给出
\[
G/(H\cap K)\cong\varphi(G)\le(G/H)\times(G/K).
\]

为证明右边的直积可解，记$A=G/H$、$B=G/K$。逐分量计算有
\[
[(a,b),(a',b')]=([a,a'],[b,b']).
\]
故$(A\times B)^{(1)}\le A^{(1)}\times B^{(1)}$。反向，取第二坐标为$e$可得到$([a,a'],e)$；取第一坐标为$e$可得到$(e,[b,b'])$，这些元素的乘积生成$A^{(1)}\times B^{(1)}$。因此两边相等。归纳得
\[
(A\times B)^{(i)}=A^{(i)}\times B^{(i)}.
\]
若$A^{(r)}=\{e\}$、$B^{(s)}=\{e\}$，取$t=\max\{r,s\}$，得到$(A\times B)^{(t)}=\{(e,e)\}$，直积可解。由\exref{2}{42}，可解群的子群$\varphi(G)$可解，同构将此性质传给$G/(H\cap K)$。
\end{proof}

\paragraph{复习提示}
$\varphi$未必满到整个直积，因此只写$G/(H\cap K)\cong\varphi(G)$，不能直接写成与整个直积同构。


\subsection{习题46：\texorpdfstring{只有两个内自同构时原群交换}{只有两个内自同构时原群交换}}
\label{ex:2:46}

\paragraph{题目}
若群$G$只有两个内自同构，证明$G$交换。进一步说明这一前提实际上不可能发生。

\paragraph{考点与思路}
利用$\operatorname{Inn}(G)\cong G/Z(G)$，将自同构数目转化为中心商群的阶数。

\paragraph{所用结论与原文依据}
\noteref{inner}；\noteref{cyclic}；\exref{2}{39}。

% 原解来源：chapters/Part 2/chapter2-exercises.tex:940；第2章原\exref{2}{46}。

\begin{proof}
假设$|\operatorname{Inn}(G)|=2$。由本章开始说明的同构，$G/Z(G)$是2阶群；2阶群必循环，故\exref{2}{39}推出$G$交换。

然而，交换群中任意$g,x$都有$gxg^{-1}=x$，所以所有内自同构都是恒等映射，必有$|\operatorname{Inn}(G)|=1$。这与假设等于2矛盾。因此不存在内自同构恰有两个的群，而题目的条件命题成立。
\end{proof}

\paragraph{复习提示}
“推出交换”后还可以继续检查前提：这里不是找到一类特殊交换群，而是证明前提本身不可能满足。


\subsection{习题47：\texorpdfstring{非平凡有限群至少有两个自同构，除了2阶群}{非平凡有限群至少有两个自同构，除了2阶群}}
\label{ex:2:47}

\paragraph{题目}
设$G$为有限群，$|G|>2$。证明$|\Aut(G)|\ge2$；2阶群是非平凡有限群中的例外。

\paragraph{考点与思路}
按是否交换、是否有阶大于2的元素分类，分别构造非恒等内自同构、取逆映射、交换基向量的线性自同构。

\paragraph{所用结论与原文依据}
\noteref{inner}；\noteref{aut}；\noteref{glaws}。补充前置工具：有限维二元域向量空间的基构造，在本题具体展开。

% 原解来源：chapters/Part 2/chapter2-exercises.tex:957；第2章原\exref{2}{47}。

\begin{proof}
恒等映射始终是自同构，因此只需再构造一个非恒等自同构。

\textbf{第一种情况：$G$不交换。}存在$a\notin Z(G)$，因而有$x\in G$使$ax\ne xa$。内自同构$c_a:y\mapsto aya^{-1}$满足$c_a(x)\ne x$，所以不是恒等映射。

\textbf{第二种情况：$G$交换，且有元素$a$的阶大于2。}取逆映射$f(x)=x^{-1}$满足
\[
f(xy)=y^{-1}x^{-1}=x^{-1}y^{-1}=f(x)f(y).
\]
又$f^2$为恒等映射，故$f$双射，是自同构。因为$a^2\ne e$，有$a^{-1}\ne a$，所以$f$非恒等。

\textbf{第三种情况：$G$交换，且所有元素平方均为$e$。}将群运算记成加法，单位元记成0，则$x+x=0$对所有$x$成立。规定$\F_2$上的标量乘法为$0x=0$、$1x=x$。向量空间的分配律中唯一需要特别检查的是
\[
(1+1)x=0=x+x;
\]
其余标量组合均直接满足，因此$G$成为$\F_2$上的向量空间。

因$G$有限，可从非零向量逐次加入不在已有张成空间中的向量，有限步后得到一组基$v_1,\ldots,v_d$。每个向量唯一写成$\sum_i\varepsilon_iv_i$，$\varepsilon_i\in\{0,1\}$，故$|G|=2^d$。由$|G|>2$得$d\ge2$。定义
\[
f\left(\varepsilon_1v_1+\varepsilon_2v_2+\sum_{i=3}^d\varepsilon_iv_i\right)
=\varepsilon_1v_2+\varepsilon_2v_1+\sum_{i=3}^d\varepsilon_iv_i.
\]
基表示唯一，故$f$定义良好；它保持加法，且$f^2$为恒等映射，因此是群自同构。又$f(v_1)=v_2\ne v_1$，所以非恒等。

三种情况覆盖所有群，均已构造第二个自同构。若$|G|=2$，唯一非单位元必须被自同构固定，单位元也被固定，因此只有恒等自同构，确为例外。
\end{proof}

\paragraph{复习提示}
取逆在一般非交换群中不是同态；先分出交换情形才能使用。所有元素平方为单位元时，取逆映射又恰是恒等，需要另作构造。


\subsection{习题48：\texorpdfstring{克莱因四元群的自同构群}{克莱因四元群的自同构群}}
\label{ex:2:48}

\paragraph{题目}
$\Aut(V_4)\cong S_3$。

\paragraph{考点与思路}
将自同构限制到三个非单位元，证明任意三点置换都保持乘法。

\paragraph{所用结论与原文依据}
\noteref{aut}；\noteref{iso}。

% 原解来源：chapters/Part 2/chapter2-exercises.tex:971；第2章原\exref{2}{48}。

\begin{proof}
记$V_4=\{e,a,b,c\}$，其中$a^2=b^2=c^2=e$、$ab=ba=c$，其余不同非单位元的乘积也等于第三个非单位元。

自同构固定$e$，且是双射，所以必置换集合$\{a,b,c\}$。定义限制映射
\[
\Phi:\Aut(V_4)\to S(\{a,b,c\}).
\]
限制保持复合，故$\Phi$为同态。若限制为恒等置换，自同构既固定三个非单位元又固定$e$，所以为恒等自同构，故$\Phi$单射。

任取$\pi\in S(\{a,b,c\})$，将它延拓为固定$e$的双射$\widetilde\pi:V_4\to V_4$。若一个因子为$e$，乘法显然保持；若两个因子为同一个非单位元，乘积为$e$，像的乘积也为$e$；若为两个不同非单位元，乘积是余下的第三个，置换后仍满足这一规则。所以$\widetilde\pi$保乘法，为自同构。因此$\Phi$满射。

由$S(\{a,b,c\})\cong S_3$，得$\Aut(V_4)\cong S_3$。
\end{proof}

\paragraph{复习提示}
一个任意的元素置换通常不保持群乘法，本题需用$V_4$的特殊乘法逐类验证。


\subsection{习题49：\texorpdfstring{$S_3,S_4$为完全群}{S3,S4为完全群}}
\label{ex:2:49}

\paragraph{题目}
证明$S_3$、$S_4$都是完全群，即中心平凡且全部自同构均为内自同构。

\paragraph{考点与思路}
先证明中心平凡以数出内自同构，再让任意自同构作用在一个小的、由群自身性质确定的集合上，给自同构总数一个同样的上界。

\paragraph{所用结论与原文依据}
\noteref{inner}；\noteref{complete}；\noteref{conjugation}；\noteref{threecycles}；\exref{2}{6}、\exref{2}{31}。

% 原解来源：chapters/Part 2/chapter2-exercises.tex:990；第2章原\exref{2}{49}。

\begin{proof}
\textbf{第一步：中心平凡。}事实上$n\ge3$时$Z(S_n)=\{e\}$。若$z\in Z(S_n)$，则对每个$i\ne j$，
\[
(z(i)\ z(j))=z(ij)z^{-1}=(ij),
\]
所以$z$保持每个二元子集$\{i,j\}$。固定$i$，取互不相同的$i,j,k$，则$z(i)$同时属于$\{i,j\}$与$\{i,k\}$，只能等于$i$。故$z=e$。

因此$\operatorname{Inn}(S_3)\cong S_3$有6个元素，$\operatorname{Inn}(S_4)\cong S_4$有24个元素。只需证明各自的全部自同构数不超过这两个数。

\textbf{第二步：$S_3$。}自同构保持元素阶，所以置换$S_3$中的三个2阶元素$(12),(13),(23)$。这给出同态
\[
\Aut(S_3)\longrightarrow\operatorname{Sym}\{(12),(13),(23)\}\cong S_3.
\]
若自同构固定这三个元素，因它们生成$S_3$，该自同构对任意生成元乘积也固定，故为恒等映射。因此此同态单射，$|\Aut(S_3)|\le6$。内自同构已经有6个，故全部自同构都是内自同构，$S_3$完全。

\textbf{第三步：$S_4$。}由\exref{2}{31}，$S_4$有四个西罗3子群。记$Q_i$为固定点$i$、在其余三个点上作3-轮换的那个3阶子群。令
\[
H_i=\{\sigma\in S_4:\sigma(i)=i\}.
\]
我们先证明$N_{S_4}(Q_i)=H_i$。若$g$正规化$Q_i$，它把$Q_i$中非单位3-轮换共轭为同一子群中的非单位3-轮换。共轭将移动点集合映为移动点集合，因此$g$保持$\{1,2,3,4\}\setminus\{i\}$，从而$g(i)=i$。反过来，若$g$固定$i$，它共轭一个移动另外三个点的3-轮换后，仍是这三个点上的3-轮换，只能为原轮换或其逆，所以正规化$Q_i$。等式得证。

自同构把西罗3子群互相置换，并满足
\[
\alpha(N_{S_4}(Q))=N_{S_4}(\alpha(Q)).
\]
故它们也把四个$H_i$互相置换，给出同态
\[
\Psi:\Aut(S_4)\longrightarrow\operatorname{Sym}\{H_1,H_2,H_3,H_4\}\cong S_4.
\]
四个$H_i$彼此不同：若$i\ne j$，取$k\ne i,j$，则$(jk)$固定$i$而不固定$j$，所以$H_i\ne H_j$。

下面证明$\Psi$单射。设$\alpha\in\ker\Psi$，即$\alpha(H_i)=H_i$对每个$i$成立。任取$g\in S_4$，共轭点稳定子得到
\[
gH_ig^{-1}=H_{g(i)}.
\]
对该式应用$\alpha$，并利用它固定各个$H_i$，有
\begin{align*}
H_{\alpha(g)(i)}
&=\alpha(g)H_i\alpha(g)^{-1}\\
&=\alpha(gH_ig^{-1})\\
&=\alpha(H_{g(i)})=H_{g(i)}.
\end{align*}
因四个$H_i$彼此不同，必有$\alpha(g)(i)=g(i)$对所有$i$成立。所以$\alpha(g)$与$g$是相同置换。$g$任意，故$\alpha$为恒等自同构，$\ker\Psi$平凡。

于是$|\Aut(S_4)|\le|S_4|=24$。结合已有24个内自同构，得到$\Aut(S_4)=\operatorname{Inn}(S_4)$。再结合第一步的中心平凡，$S_4$也是完全群。
\end{proof}

\paragraph{复习提示}
证明自同构作用的核平凡是本题关键。仅说“自同构置换四个西罗3子群”，还不能直接断言自同构至多24个；必须排除作用不忠实的可能。


\subsection{习题50：\texorpdfstring{非交换单群的自同构群完全}{非交换单群的自同构群完全}}
\label{ex:2:50}

\paragraph{题目}
若$G$为非交换单群，则$\Aut(G)$是完全群。

\paragraph{考点与思路}
先证明内自同构群是唯一极小非平凡正规子群，再把自同构在它上的限制消去。

\paragraph{所用结论与原文依据}
\noteref{inner}；\noteref{autcenterless}；\noteref{characteristic}；\noteref{complete}；\noteref{simple}。

% 原解来源：chapters/Part 2/chapter2-exercises.tex:1020；第2章原\exref{2}{50}。

\begin{proof}
设$A=\Aut(G)$、$S=\Inn(G)$。因为$G$非交换且单，中心作为正规子群只能为$\{e\}$：若中心为$G$，群便交换。因此原笔记的内自同构群定理给出$S\cong G$，所以$S$也是非交换单群，且$S\trianglelefteq A$。

由原笔记无中心群的自同构定理，$C_A(S)=\{e_A\}$且$Z(A)=\{e_A\}$。这里$C_A(S)$指$A$中与$S$每个元素交换的自同构所成的子群。

先证明$S$是$A$的特征子群。任取非平凡正规子群$M\trianglelefteq A$。交$M\cap S$为$S$的正规子群，由$S$的单性，只能平凡或等于$S$。若交平凡，因$M,S$均在$A$中正规，对$m\in M,s\in S$，换位子$[m,s]$同时落在$M,S$中，故为单位元。于是$M\leq C_A(S)=\{e_A\}$，矛盾。故每个非平凡正规子群$M$都包含$S$。

$S$自身没有非平凡真正规子群，所以它是$A$中唯一的极小非平凡正规子群。任意$A$的自同构都把极小非平凡正规子群映为同类子群，因此必须保持$S$，故$S$为特征子群。

任取$\beta\in\Aut(A)$。因$S$特征，$\beta|_S$为$S$的自同构。通过同构$G\to S$、$g\mapsto c_g$（$c_g(x)=gxg^{-1}$），该限制对应某个$a\in\Aut(G)=A$。在$A$中有
\[
ac_ga^{-1}=c_{a(g)},
\]
所以由$a$诱导的内自同构$\gamma_a$与$\beta$在$S$上相同。令$\delta=\gamma_a^{-1}\beta$，则$\delta$逐元素固定$S$。

对$x\in A,s\in S$，因$S\trianglelefteq A$，$xsx^{-1}\in S$，从而
\[
\delta(x)s\delta(x)^{-1}=\delta(xsx^{-1})=xsx^{-1}.
\]
因此$x^{-1}\delta(x)$与所有$s\in S$交换，属于$C_A(S)=\{e_A\}$，故$\delta(x)=x$。所以$\delta$恒等，$\beta=\gamma_a$为内自同构。结合$Z(A)$平凡，按定义$A$为完全群。
\end{proof}

\paragraph{复习提示}
关键是先证明$S$特征，才能把$A$的任意自同构限制为$S$的自同构。


\subsection{习题51：\texorpdfstring{不存在换位子群为$S_4$的更大群}{不存在换位子群为S4的更大群}}
\label{ex:2:51}

\paragraph{题目}
不存在$G\supsetneq S_4$使$G^{(1)}=S_4$。

\paragraph{考点与思路}
沿原笔记用$S_4$的完全性把整个群分解成它与中心化子的直积，再取换位子群。

\paragraph{所用结论与原文依据}
\noteref{complete}；\noteref{characteristic}；\noteref{internal}；\noteref{derived}；\noteref{threecycles}；\noteref{sign}；\exref{2}{49}。

% 原解来源：chapters/Part 2/chapter2-exercises.tex:1052；第2章原\exref{2}{51}。

\begin{proof}
反设$N=S_4=G^{(1)}$。换位子群正规，所以$N\trianglelefteq G$；由第二章习题49，$N$中心平凡且为完全群。

任取$g\in G$。由于$N$正规，$x\mapsto gxg^{-1}$限制到$N$为自同构。$N$完全，故存在$n\in N$使$gxg^{-1}=nxn^{-1}$对所有$x\in N$成立。因此$n^{-1}g\in C_G(N)$，即$g\in NC_G(N)$。于是$G=NC_G(N)$。

记$C=C_G(N)$。由定义$N$与$C$逐元素交换，且$N\cap C=Z(N)=\{e\}$。故乘法映射$N\times C\to G$、$(n,c)\mapsto nc$是满同态；若$nc=e$，则$n=c^{-1}\in N\cap C$，所以$n=c=e$，核平凡。因而$G\cong N\times C$。

对直积逐分量计算换位子，并分别用$[(n_1,e),(n_2,e)]$、$[(e,c_1),(e,c_2)]$取得两因子的生成换位子，得到
\[
G^{(1)}=N^{(1)}C^{(1)}.
\]
因为$C^{(1)}\leq G^{(1)}=N$且$C^{(1)}\leq C$，得$C^{(1)}\leq N\cap C=\{e\}$。所以$G^{(1)}=N^{(1)}$。

最后说明$N^{(1)}=A_4$：符号同态使所有换位子落入$A_4$；另一方面两个对换的换位子可给出3-轮换，例如$[(12),(23)]=(132)$，对符号作共轭替换即可得到全部3-轮换。原笔记已证明3-轮换生成$A_4$，故$N^{(1)}=A_4$。于是$G^{(1)}=A_4\ne S_4$，与假设矛盾。
\end{proof}

\paragraph{复习提示}
直积分解后还要证明中心化子因子的导群平凡，不能直接省略这一因子。


\subsection{习题52：\texorpdfstring{3-轮换为$A_n$中的换位子}{3-轮换为An中的换位子}}
\label{ex:2:52}

\paragraph{题目}
$n>4$时每个3-轮换为一个换位子，从而$A_n^{(1)}=A_n$。

\paragraph{考点与思路}
保留原解的两个双对换，直接计算它们的换位子。

\paragraph{所用结论与原文依据}
\noteref{derived}；\noteref{threecycles}；\noteref{cycles}。

% 原解来源：chapters/Part 2/chapter2-exercises.tex:1073；第2章原\exref{2}{52}。

\begin{proof}
任取不同的$a,b,c$。因$n\geq5$，还可选取与它们不同且彼此不同的$d,e$。令
\[
x=(ab)(de),\qquad y=(bc)(de).
\]
两个置换均为两个对换之积，故$x,y\in A_n$。它们又满足$x^{-1}=x$、$y^{-1}=y$。不相交支集上的置换交换，所以
\[
xy=(ab)(bc)(de)^2=(abc).
\]
按本笔记的换位子约定，
\[
[x,y]=x^{-1}y^{-1}xy=(xy)^2=(acb),\qquad
[y,x]=(abc).
\]
因此每个3-轮换都是$A_n$中某两个元素的换位子，属于$A_n^{(1)}$。

全部3-轮换生成$A_n$，所以$A_n\leq A_n^{(1)}$。反向包含由换位子群是原群的子群而成立。因此$A_n^{(1)}=A_n$。
\end{proof}

\paragraph{复习提示}
附加的两个点保证构造的$x,y$都是偶置换，这解释了$n>4$的用途。


\subsection{习题53：\texorpdfstring{$S_4$的两个生成元及关系}{S4的两个生成元及关系}}
\label{ex:2:53}

\paragraph{题目}
存在$a,b$生成$S_4$，满足$a^2=b^3=(ab)^4=e$。

\paragraph{考点与思路}
按原解取一个对换和一个3-轮换，用共轭取得星形对换。

\paragraph{所用结论与原文依据}
\noteref{conjugation}；\noteref{cycles}；\noteref{generated}；\exref{2}{10}。

% 原解来源：chapters/Part 2/chapter2-exercises.tex:1092；第2章原\exref{2}{53}。

\begin{proof}
取$a=(12)$、$b=(234)$。它们分别为2阶和3阶，故$a^2=b^3=e$。逐点作用可知$ab$把$1,2,3,4$依次映到$2,3,4,1$，所以$ab=(1234)$，从而$(ab)^4=e$。

设$K=\langle a,b\rangle$。\proofref{conjugation}{共轭公式}给出
\[
bab^{-1}=(13),\qquad b^2ab^{-2}=(14).
\]
因此$K$含$(12),(13),(14)$。由第二章习题10，这三个星形对换生成$S_4$，所以$S_4\leq K$；反向包含成立，故$K=S_4$。

由此构造出了满足要求的生成元。这里证明的是这两个具体置换生成$S_4$，并未把题中的三个等式单独视为任意群的完整展示。
\end{proof}

\paragraph{复习提示}
验证所列等式和证明两个元素生成全群，是两个分别需要完成的步骤。


\subsection{习题54：\texorpdfstring{$A_4$的两个生成元及关系}{A4的两个生成元及关系}}
\label{ex:2:54}

\paragraph{题目}
存在$a,b$生成$A_4$，满足$a^2=b^3=(ab)^3=e$。

\paragraph{考点与思路}
先计算关系，再让三个共轭双对换构成$V_4$并用子群乘积计数。

\paragraph{所用结论与原文依据}
\noteref{conjugation}；\noteref{sign}；\noteref{lagrange}；\exref{1}{11}、\exref{1}{21}。

% 原解来源：chapters/Part 2/chapter2-exercises.tex:1110；第2章原\exref{2}{54}。

\begin{proof}
取$a=(12)(34)$、$b=(123)$。显然$a^2=b^3=e$。按右边先作用计算，$ab$固定1，并将$2,4,3$依次映到$4,3,2$，所以$ab=(243)$，从而$(ab)^3=e$。

设$K=\langle a,b\rangle$。$a,b$均为偶置换，所以$K\leq A_4$。由\proofref{conjugation}{共轭公式}，
\[
bab^{-1}=(23)(14),\qquad b^2ab^{-2}=(13)(24).
\]
这两者连同$a$恰为$V_4$的三个非单位元，所以$V_4\leq K$。又$\langle b\rangle$阶为3，$V_4$阶为4，两者的交的阶同时整除3和4，故交平凡。

由第一章习题21，集合$V_4\langle b\rangle$有$4\cdot3=12$个元素，并且包含于$K$。而$K\leq A_4$、$|A_4|=12$，所以只能$K=A_4$。
\end{proof}

\paragraph{复习提示}
不必预先证明$V_4\langle b\rangle$为子群；集合乘积的计数已足够比较阶。


\subsection{习题55：\texorpdfstring{四元数群的关系与条件}{四元数群的关系与条件}}
\label{ex:2:55}

\paragraph{题目}
设$G=\langle a,b\rangle$，满足
\[
a^4=b^4=e,\qquad a^2=b^2,\qquad b^{-1}ab=a^{-1}.
\]
在$a^2\ne e$的条件下，证明$G\cong Q_8$，列出全部子群并证明每个子群正规。若将上述关系理解为完整展示，也证明展示群为$Q_8$。仅要求等式成立时，说明为何须加限定。

\paragraph{考点与思路}
先说明生成关系的读法，再用八个正规形识别四元数群，并按元素阶列尽子群。

\paragraph{所用结论与原文依据}
\noteref{dihedral}；\noteref{quaternion}；\noteref{lagrange}；\noteref{free}；\exref{1}{12}。

% 原解来源：chapters/Part 2/chapter2-exercises.tex:1126；第2章原\exref{2}{55}。

\begin{proof}
若只要求题设等式成立而不作其他限定，$V_4$中两个不同非单位元$a,b$也满足这些等式，却不与$Q_8$同构。因此先在$a^2\ne e$的条件下证明，再说明完整展示的情形。

由$b^{-1}ab=a^{-1}$得$ab=ba^{-1}$；将$a$换成$a^{-1}$又得$ba=a^{-1}b$。利用这些关系可以把每个字中的$b$移到右边，并用$b^2=a^2$、$a^4=e$化简，故每个元素都能写成
\[
a^ib^j,\qquad0\leq i<4,\quad0\leq j<2.
\]
所以$|G|\leq8$。$a^2\ne e$与$a^4=e$说明$a$阶为4。若$b\in\langle a\rangle$，则$b$与$a$交换，由共轭关系得$a=a^{-1}$，与4阶矛盾。因此$b\notin\langle a\rangle$，两个右陪集$\langle a\rangle$、$\langle a\rangle b$不交，每个均有四个元素，故$|G|=8$，所列正规形两两不同。

在四元数体中，$Q_8=\{\pm1,\pm I,\pm J,\pm K\}$是乘法群，其中$I^2=J^2=K^2=-1$、$IJ=K=-JI$。元素$I,J$满足与$a,b$相同的关系。将$a^ib^j$映到$I^iJ^j$，两边八个正规形唯一，故为双射。化简乘积时均使用$ba=a^{-1}b$或$JI=I^{-1}J$及平方关系，所得乘法一致，所以映射为同构，$G\cong Q_8$。

若将题设关系理解为完整展示，设相应展示群为$\widetilde G$。同样的化简证明$|\widetilde G|\leq8$。$I,J$满足全部关系，由自由群商构造的泛性质，$a\mapsto I,b\mapsto J$诱导$\widetilde G\twoheadrightarrow Q_8$，所以$|\widetilde G|\geq8$，从而也与$Q_8$同构。

最后列出全部子群。$Q_8$只有一个2阶元素$-1$，六个其余非单位元均为4阶，它们分别生成
\[
\langle I\rangle=\{1,-1,I,-I\},\quad
\langle J\rangle=\{1,-1,J,-J\},\quad
\langle K\rangle=\{1,-1,K,-K\}.
\]
由\proofref{lagrange}{拉格朗日定理}，真子群的阶只能为1、2或4。1阶子群只有$\{1\}$；2阶子群只能为$\{1,-1\}$；4阶子群不可能所有非单位元都为2阶，因为全群只有一个2阶元，故必须含4阶元，只能为上述三个循环子群之一。因此连同全群$Q_8$，已穷尽全部子群。

三个4阶子群指数为2，均正规；$\{1,-1\}$包含于中心，也正规；平凡子群及全群本来就正规。所以$Q_8$的每个子群均正规。
\end{proof}

\paragraph{复习提示}
仅有若干成立的等式，与把它们指定为完整生成关系，是两种不同假设。


\subsection{习题56：\texorpdfstring{60以下的群及60阶单群}{60以下的群及60阶单群}}
\label{ex:2:56}

\paragraph{题目}
每个阶小于60的群都可解；其中单群只可能为素数阶循环群。60阶单群同构于$A_5$。

\paragraph{考点与思路}
保留原笔记的六步路线：用单群的置换作用排除小阶情形，以归纳证明可解性，再从60阶单群的西罗子群作用识别$A_5$。

\paragraph{所用结论与原文依据}
\noteref{simple}；\noteref{action}；\noteref{firstiso}；\noteref{sylow1}；\noteref{sylow2}；\noteref{sylow3}；\noteref{pcenter}；\noteref{solvable}；\noteref{an}；\exref{2}{20}、\exref{2}{32}。补充前置工具：有限阶数分类与关于群阶的数学归纳法。

% 原解来源：chapters/Part 2/chapter2-exercises.tex:1161；第2章原\exref{2}{56}。

\begin{proof}
\textbf{第一步：建立单群的置换作用限制。}
非交换单群$G$的非平凡作用$\rho:G\to S_t$有正规核。
核不能为$G$，所以只能平凡，$\rho$忠实。
符号复合同态$\operatorname{sgn}\circ\rho$若不平凡，同样须单射到$C_2$，与非交换性矛盾。
故像包含于$A_t$，得到
\[
|G|\mid t!/2.
\]
特别地，西罗子群的共轭作用在单群中必须满足$n_p>1$，并要求$|G|\mid n_p!/2$。

\textbf{第二步：排除60以下的非交换单群。}
素数幂阶群有非平凡中心，若单就只能交换，因而只可能为素数阶循环群。
对于其余合数$n<60$，逐项使用$n_p\mid n/p^{v_p(n)}$及$n_p\equiv1\pmod p$。
其中$v_p(n)$表示$n$中素因子$p$的指数。若西罗$p$子群正规且非平凡、真包含于群，则与单性矛盾。
以下阶数选择其最大素因子$p$就迫使$n_p=1$：
\[
\begin{gathered}
6,10,14,15,18,20,21,22,26,28,33,34,35,38,39,\\
40,42,44,45,46,50,51,52,54,55,57,58.
\end{gathered}
\]

\begin{center}\begin{tabular}{c|p{11cm}}
$p$ & 最大素因子为$p$且西罗数被迫为1的阶数$n$\\\hline
$3$ & $6,18,54$\\
$5$ & $10,15,20,40,45,50$\\
$7$ & $14,21,28,35,42$\\
$11$ & $22,33,44,55$\\
$13$ & $26,39,52$\\
$17$ & $34,51$\\
$19$ & $38,57$\\
$23$ & $46$\\
$29$ & $58$\\
\end{tabular}\end{center}
例如$n=45$时$n_5\mid9$，可选1、3、9中只有1模5余1；其余同样由有限个因子核对。
余下仅有$12,24,30,36,48,56$。
12、24、48阶若不已有正规西罗2子群，则$n_2=3$，第一步要求嵌入$A_3$，不可能；36阶若$n_3\neq1$，就有$n_3=4$，要求嵌入$A_4$，也不可能。
30阶单群须$n_5=6,n_3=10$，不同素数阶子群贡献互不重叠的非单位元，至少有
\[
6(5-1)+10(3-1)=44>29,
\]
矛盾。56阶由\exref{2}{32}排除。
所以60以下单群只有素数阶循环群。

\textbf{第三步：证明60以下的全部群可解。}
对阶归纳，平凡群可解。
若$G$单，第二步说明它为素数阶循环群，因而可解。
若不单，取非平凡真正规子群$N$，$N$及$G/N$阶都更小，归纳得二者可解，再由扩张性质得$G$可解。

\textbf{第四步：计算60阶单群的西罗数及元素数。}
设$|G|=60$且$G$单，则由计数及第一步，
\[
n_5=6,\qquad n_3=10.
\]
这里$n_5$只能为1或6；$n_3$只能为1、4、10，1由单性排除，4由$60\nmid|A_4|$排除。
因此5阶元有24个，3阶元有20个，余下非单位元只有
\[
60-1-24-20=15
\]
个。西罗2子群阶为4，其个数只能为1、3、5、15，前两者由单性及$A_3$作用限制排除。

\textbf{第五步：排除$n_2=15$。}
若西罗2子群循环，则15个共轭子群全循环，每个给出两个4阶元；4阶元唯一生成它所属的4阶子群，故不同子群的这些元素不重叠。
于是至少有30个4阶元，超过剩余15个，矛盾。
所以这些西罗子群均为$V_4$，每个含三个对合（即2阶元素），合计45次出现。
全部对合至多15个，所以至少有一个对合$t$属于两个不同的西罗2子群。
这两个$V_4$中的元素都与$t$交换，故它们都包含于$C_G(t)$，使
\[
4\mid|C_G(t)|,\qquad|C_G(t)|>4.
\]
由其阶整除60，只能为12、20或60。
60使$t$属于非交换单群的中心，不可能；20给出指数3的陪集作用，同样由第一步排除。
若为12，则指数5的陪集作用使$G\hookrightarrow A_5$，两群阶相等，故$G\cong A_5$。
但$A_5$中不存在4阶元素：5个点上的4阶置换只能含一个4-轮换，而4-轮换为奇置换。因此它的4阶西罗2子群都为$V_4$。这些子群只有五个：其15个对合均为双对换，每个固定唯一一点，固定给定一点的三个双对换加$e$构成一个$V_4$。
对$t=(12)(34)$，$S_5$中与它交换的置换须固定5，可交换两个二元组并在组内交换，所以中心化子阶为$2^2\cdot2!=8$。
其中含奇置换$(12)$，故偶置换部分恰占一半：
\[
|C_{A_5}(t)|=4.
\]
它已经包含固定5的那个$V_4$，所以恰等于它；其他双对换由共轭有相同结论。
若某$V_4$含$t$，就必须包含于这个4阶中心化子，因而等于它。
所以五个这样的子群已列尽，与$n_2=15$矛盾。

\textbf{第六步：用西罗2子群的共轭作用识别群。}
只能$n_2=5$。$G$共轭作用于这五个子群，给出非平凡同态$G\to S_5$。
单性及符号同态的讨论使它单射到$A_5$，再比较阶，得到
\[
G\cong A_5.
\]
\end{proof}

\paragraph{复习提示}
在假设单群时作用核只有两种可能，先证明作用非平凡才能推出忠实；计数时还要防止不同4阶子群的2阶元重复。


\subsection{习题57：\texorpdfstring{有限指数子群仍有限生成}{有限指数子群仍有限生成}}
\label{ex:2:57}

\paragraph{题目}
若$G$有限生成且$[G:H]=m<\infty$，则$H$有限生成。

\paragraph{考点与思路}
选定有限左陪集代表，从一个字的右端递推代表元，得到有限个施赖尔生成元。

\paragraph{所用结论与原文依据}
\noteref{generated}；\noteref{cosets}。

% 原解来源：chapters/Part 2/chapter2-exercises.tex:1243；第2章原\exref{2}{57}。

\begin{proof}
设$G=\langle g_1,\ldots,g_n\rangle$。令$S=\{g_1^{\pm1},\ldots,g_n^{\pm1}\}$，它有限，且每个群元素都是$S$中元素的有限乘积。取$H$的$m$个左陪集代表组成集合$T$，规定$e$代表陪集$H$。对$x\in G$，用$\overline x\in T$表示$xH$的唯一选定代表。

对每个$s\in S,t\in T$，由$stH=\overline{st}H$知
\[
b_{s,t}=\overline{st}^{\,-1}st\in H.
\]
令$\mathcal B=\{b_{s,t}:s\in S,t\in T\}$，则$\mathcal B\subseteq H$且$|\mathcal B|\leq|S|m$，所以$\langle\mathcal B\rangle\leq H$。

任取$h\in H$，写$h=s_1\cdots s_r$，$s_i\in S$。$r=0$时$h=e$已属于$\langle\mathcal B\rangle$。以下设$r\geq1$，从右端递归定义
\[
t_{r+1}=e,\qquad t_i=\overline{s_it_{i+1}}quad(i=r,r-1,\ldots,1).
\]
由反向归纳可知$t_iH=s_i\cdots s_rH$：起点$t_{r+1}H=H$；若$t_{i+1}H=s_{i+1}\cdots s_rH$，则
\[
t_iH=s_it_{i+1}H=s_i\cdots s_rH.
\]
因为$h\in H$，有$t_1H=H$，按代表元的选取，$t_1=e$。

令$b_i=t_i^{-1}s_it_{i+1}$。由$t_i=\overline{s_it_{i+1}}$，每个$b_i$均为$\mathcal B$的元素。相邻代表元消去，得到
\begin{align*}
b_1\cdots b_r
&=t_1^{-1}s_1t_2\,t_2^{-1}s_2t_3\cdots t_r^{-1}s_rt_{r+1}\\
&=t_1^{-1}s_1\cdots s_rt_{r+1}=h.
\end{align*}
故$H\leq\langle\mathcal B\rangle$，结合反向包含，得$H=\langle\mathcal B\rangle$，所以$H$有限生成。
\end{proof}

\paragraph{复习提示}
写成$a_j^{-1}g_i a_k$时，$a_j$必须代表$g_i a_kH$，不能让三个下标完全独立。


